\documentclass[12pt,a4paper]{article}

% Encoding and fonts
\usepackage[utf8]{inputenc}
\usepackage[T1]{fontenc}
\usepackage{lmodern}

% Layout
\usepackage[margin=1in]{geometry}
\usepackage{parskip}
\usepackage{calc}

% Links and references
\usepackage[colorlinks=true,linkcolor=blue!70!black,urlcolor=blue!70!black,citecolor=blue!70!black,breaklinks=true]{hyperref}
\usepackage{bookmark}

% Code listings with Go syntax highlighting
\usepackage{listings}
\usepackage{xcolor}

\definecolor{goKeyword}{RGB}{0,0,180}
\definecolor{goComment}{RGB}{106,115,125}
\definecolor{goString}{RGB}{163,21,21}
\definecolor{goType}{RGB}{38,127,153}
\definecolor{goBackground}{RGB}{248,248,248}
\definecolor{goLineNum}{RGB}{160,160,160}

\lstdefinelanguage{Gno}{
  language=Go,
  morekeywords={realm,cross,cur,attach,safely,revive},
  morekeywords=[2]{string,int,bool,any,error,byte,rune,uint64},
  morekeywords=[3]{true,false,nil,iota},
  sensitive=true,
}

\lstset{
  language=Gno,
  basicstyle=\ttfamily\footnotesize\linespread{0.9}\selectfont,
  keywordstyle=\color{goKeyword}\bfseries,
  keywordstyle=[2]\color{goType},
  keywordstyle=[3]\color{goKeyword},
  commentstyle=\color{goComment},
  stringstyle=\color{goString},
  breaklines=true,
  frame=single,
  backgroundcolor=\color{goBackground},
  showstringspaces=false,
  tabsize=4,
  numbers=left,
  numberstyle=\tiny\color{goLineNum},
  numbersep=8pt,
  xleftmargin=2em,
  framexleftmargin=1.5em,
  captionpos=b,
  aboveskip=0.8em,
  belowskip=0.4em,
  rulecolor=\color{gray!40},
  columns=fullflexible,
  keepspaces=true,
}

% Tables and lists
\usepackage{longtable}
\usepackage{array}
\usepackage{enumitem}

% Graphics
\usepackage{graphicx}

% Allow deep nesting of itemize (up to 9 levels)
\setlistdepth{9}
\renewlist{itemize}{itemize}{9}
\setlist[itemize,1]{label=\textbullet}
\setlist[itemize,2]{label=--}
\setlist[itemize,3]{label=*}
\setlist[itemize,4]{label=\textperiodcentered}
\setlist[itemize,5]{label=$\diamond$}
\setlist[itemize,6]{label=$\circ$}
\setlist[itemize,7]{label=\textbullet}
\setlist[itemize,8]{label=--}
\setlist[itemize,9]{label=*}

% Pandoc compatibility
\providecommand{\tightlist}{\setlength{\itemsep}{0pt}\setlength{\parskip}{0pt}}
\newcommand{\passthrough}[1]{#1}

\title{\Huge\textbf{Gno.land} \\ \Large Manifesto \& Interrealm Specification}
\author{Jae Kwon}
\date{March 9, 2026}

\begin{document}

\maketitle
\tableofcontents
\newpage

\part{Gno.land Manifesto}
\label{gno.land-manifesto}
\setcounter{section}{1}

\emph{@author: Jae Kwon}

\emph{The prefix ``gno'' in Koine Greek is derived from the verb
``ginōskō'', which means ``to know'' or ``to recognize.'' It is often
associated with terms related to knowledge, such as ``gnosis'', which
signifies knowledge or insight, particularly in a spiritual context.}

To jump to the technical portions see \hyperref[gno-language]{Gno
Language} and \hyperref[gno-land-blockchain]{Gno.land Blockchain}.

% Table of Contents subsection removed (LaTeX generates its own)

\subsection{Introduction}\label{introduction}

\subsubsection{Gno.land Raison d'Etre}\label{gno.land-raison-detre}

Gno.land was inspired by the desire to address intentional
mistranslations of the bible. At the moment of writing chances are that
the bible you have or find contains mistranslations from the original
Hebrew or Koine Greek meant to keep you enslaved to authority. Coding
for Gno.land began during the Covid19 lockdowns when I was most
frustrated at the censorship of information regarding the true
laboratory origins of Covid19 (created in the Wuhan lab as attested by
multiple whistleblowers and researchers and made possible by Fauci) and
the conspiracy to forcibly medicate the population with a dangerous
experimental gene therapy that did more harm to children and young
adults than it did good nor prevented transmission (and this was known
at the time of the mandates).

The censorship of alternative narratives during this time was a wake-up
call that we need to invert our information architecture to be open,
transparent, and \emph{forkable} instead of closed and centralized. This
follows from a fundamental mismatch in incentives: the business ``moat''
of centralized online services is generally all the data that they
``own'' contributed voluntarily by the users at first, and reluctantly
later as network effects make exit near impossible; and this in turn in
combination with profit incentives makes the centralized abuse of the
``moat'' the norm rather than exception.

The only way to overcome the ``satanification'' of our devices and the
internet, the bedrock of the information age, is to invert control by
being \emph{uncompromising} in the Seven Pillars of Good Systems:

\begin{itemize}
\tightlist
\item
  FOSS (free \& open source) software \& hardware
\item
  security (simplicity, audits, formal verification, absolute Kelvin
  versioning)
\item
  accountability (e.g.~cryptographic signatures, economic bonds, commits
  and proofs)
\item
  modularity (do one thing well; interoperable components for faster
  innovation)
\item
  transparency (if possible and reasonable)
\item
  decentralization (of validation and data silos)
\item
  exitability (e.g.~forkability; permissionless competition)
\end{itemize}

It won't be an easy task to re-invent the myriad of products and
services that we have become accustomed to, but it is feasible and
necessary, and therefore also inevitable.

\subsubsection{Gno.land, the Gno language, and the Gno
VM.}\label{gno.land-the-gno-language-and-the-gno-vm.}

The Gno virtual machine and language was created to address one of the
key problems of blockchain smart-contract systems: the lack of a good
familiar modern-day general purpose programming language for Ethereum
Solidity-like massive multi-user dapp development where integration with
external dapp logic as dependencies is native to the language itself.
You can develop dapps in Go or Rust for WASM modules today but this
approach is limited--each module/contract runs in its own memory silo
with no globally shared heap space, and support for the sharing of
pointers (to objects or closures) across modules is limited. Furthermore
unlike Solidity where contract data is automatically persisted to disk
at the end of a transaction, developing in Go or Rust contracts for WASM
inherits all the burdens of data serialization and persistence (and so
of ORMs or alternatives) making smart contract development in Go or Rust
just as complex as server-side Web2 programming if not more so.

Gno.land solves the problem by virtue of a new VM design that interprets
Gno. The Gno language is Go but extended minimally to make Go aware of
different user agent contexts in execution and user storage realm
contexts in persistence for security; and the Gno VM keeps track of
creations, updates, and deletions of runtime objects at user transaction
boundaries and throughout each transaction for automatic persistence
that also computes Merkle hashes according to the shape of the object
graph. In other words, the Gno VM simulates a machine that unifies disk
and memory (i.e.~simulates memristors which will become more dominant in
the future with progress in AI).

The major benefit of dapp development in Gno therefore is the ease of
programming where multiple layers of complexity are all but
eliminated--the interoperability of contract applications (function
calls to external import realm functions), and the serialization and
persistence of data (or any mapping between memory objects and
persistent data). As compared to Ethereum Solidity, Gno programming is
not only familiar to Go developers also with ports of many of the great
tooling that developers love in Go; Gno also lets you write complex
applications in significantly less time and lines of code.

As when Ethereum was first released implemented in Python, at first
Gno.land may not be able to compete in transaction throughput as
compared to Solana or other smart contract blockchains optimized for
speed. However as the ecosystem of Gno applications and libraries grow,
so too will the set of alternative optimized implementation of the
GnoVM, and with the horizontal scaling that will be offered in
partnership with the AtomOne chain the Gno.land ecosystem will over time
become a major competitor among existing and new smart contract systems.
Furthermore the GnoVM and Gno's unique features open the door to
applications outside of the blockchain context even for the home that
follow the Seven Pillars of Good Systems.

\subsection{Problem: Misinformation and Loss of
Trust}\label{problem-misinformation-and-loss-of-trust}

The term ``fake news'' came into mainstream consciousness during the
2016 US presidential election. The Cambridge Analytica scandal made
headlines regarding the unauthorized collection of personal data from
millions of Facebook profiles:

\begin{quote}
``The data analytics firm that worked with Donald Trump's election team
and the winning Brexit campaign harvested millions of Facebook profiles
of US voters, in one of the tech giant's biggest ever data breaches, and
used them to build a powerful software program to predict and influence
choices at the ballot box.

A whistleblower has revealed to the Observer how Cambridge Analytica --
a company owned by the hedge fund billionaire Robert Mercer, and headed
at the time by Trump's key adviser Steve Bannon -- used personal
information taken without authorisation in early 2014 to build a system
that could profile individual US voters, in order to target them with
personalised political advertisements.'' -
\href{https://www.theguardian.com/news/2018/mar/17/cambridge-analytica-facebook-influence-us-election}{Carole
Cadwalladr \& Emma Graham-Harrison, The Guardian; March 17th, 2018}
\end{quote}

\begin{quote}
It was the scandal which finally exposed the dark side of the big data
economy underpinning the internet. The inside story of how one company,
Cambridge Analytica, misused intimate personal Facebook data to
micro-target and manipulate swing voters in the US election, is
compellingly told in ``The Great Hack'' -
\href{https://www.amnesty.org/en/latest/news/2019/07/the-great-hack-facebook-cambridge-analytica/}{Amnesty
International}
\end{quote}

\begin{quote}
``One of the most urgent and uncomfortable questions raised in The Great
Hack is: to what extent are we susceptible to such behavioural
manipulation?'' - Joe Westby
\end{quote}

Trust in our institutions is at an all-time low. CIA director William
Casey was accused of saying that ``we'll know our disinformation program
is complete when everything the American public believes is false.'' by
Barbara Honegger, assistant to the chief domestic policy advisor to
President Reagan. The
\href{https://truthstreammedia.com/2015/01/13/cia-flashback-well-know-our-disinformation-program-is-complete-when-everything-the-american-public-believes-is-false/}{quote
is contested} but we know now about illegal mass mind control projects
like MKUltra (by the CIA), and also that the Smith-Mundt Modernization
Act of 2013 lifted the ban on domestic dissemination of propaganda in
the US, reversing the original Smith-Mundt Act of 1948. Whether the
quote is true or not, the substance is clearly true today.

\subsubsection{Health Sovereignty}\label{health-sovereignty}

The WHO Pandemic Agreement was adopted by the World Health Assembly in
May 2025 despite the non-participation of the United States and the
abstention of 11 other countries citing sovereignty concerns. Although
the agreement claims that nothing in it grants the WHO authority to
``direct, order, alter or otherwise prescribe'' national law or to
``mandate or impose'' specific actions such as vaccination mandates or
lockdowns, the very existence of binding international frameworks for
pandemic response centralizes authority in an institution that has
already demonstrated willingness to suppress dissenting scientific views
during Covid19. Furthermore the agreement's pathogen access and
benefit-sharing provisions create new obligations around technology
transfer and intellectual property that could be weaponized against
nations that do not comply.

Given the track record of institutional capture by pharmaceutical
interests and the censorship of legitimate scientific debate during the
Covid19 pandemic, any centralized health governance framework--no matter
how carefully worded--poses an existential risk to medical freedom and
informed consent. Decentralized systems for publishing, verifying, and
discussing scientific research are therefore not merely useful but
necessary to preserve health sovereignty. Gno contracts can host
transparent peer-review systems, accountable clinical trial registries,
and censorship-resistant repositories of medical literature that no
single authority can suppress or alter.

\subsection{Problem: AI Armageddon}\label{problem-ai-armageddon}

By \href{https://www.youtube.com/watch?v=wDBy2bUICQY}{the latest
projections (``There is no AI Bubble'')} 2027 AI will be able to
accomplish by human standards week long tasks on the order of a day. By
2029 year long tasks will be possible in the same amount of time.

This means that by 2027 a single AI instance may be able to find and
execute exploits of vulnerable systems much faster than humans can
react. By the time a service administrator detects an issue and finds a
way to patch the hack attempt, the AI may have already responded with
another backdoor. And by 2029 any single vibe user will be able to plan
and execute the exploit of multiple systems, or find a new zero day
vulnerability in just a few hours. It's not possible to patch all the
systems by 2029 even if we restrict ourselves to known vulnerabilities,
let alone dealing with new zero day exploits. Realistically by 2028 it
is inevitable that we see a major crisis from an AI swarm that attacks
and continuously attacks our internet infrastructure to the point where
we cannot even understand the truth of what is happening at a global
scale, let alone recover from it.

It's not just that AI is getting more intelligent; large corporations
are building AI and robotics systems seemingly in order to bring about
armageddon. An
\href{./images/manifesto/project_lazarus_leak.jpg}{anonymous
whistleblower leaked in 2023} that Meta is building an AI that can take
over a deceased person's social media account and continue posting
convincing posts including age progression. On 5 February 2025,
Alphabet, which owns Google,
\href{https://www.stopkillerrobots.org/news/alphabet-rollback-on-policy-to-not-use-ai-for-weapons/}{reneged
on its pledge to not use AI for weapons}. Eric Schmidt former CEO of
Google now runs an AI drone company and is a licensed arms dealer.
Boston Dynamics develops the Atlas humanoid robot that can be even
\href{https://www.youtube.com/watch?v=tjFHRVr7aNE}{more agile} than
humans, and the US military works with many robotics companies such as
with Foundation Future Industries Inc.~which develops the Phantom MK1
designed for military applications including carrying firearms.

Here are some of the things to look forward to in the coming AI
armageddon.

\begin{itemize}
\tightlist
\item
  internet discourse taken over by AI
\item
  fake news indistinguishable from real news
\item
  certificate authority system failure
\item
  mass internet \& power outages
\item
  decentralized radio communications jamming
\item
  complete collapse of trust
\item
  mass hysteria
\item
  military takeover for continuity of government
\item
  hacked robots coordinated by malicious AGI
\item
  inadvertent nuclear war
\end{itemize}

We have largely speaking two paths to follow. The best path is to enact
regulations to curtail the AI arms race. Ex-Google design ethicist and
AI expert Tristan Harris \href{https://youtu.be/tphiJE01qxw?t=400}{makes
a convincing case for AI regulation}:

\begin{quote}
``If we stop or slow down, then China will build it.'' \ldots{} The
``it'' that they would keep building is the same uncontrollable AI that
we would build. So I don't see a way out of this without there being
some agreement or negotiation between powers and countries. \ldots{}
We've done this before {[}with fluorocarbons and non-proliferation of
arms{]}.
\end{quote}

However as also noted by Tristan it would be the most difficult
coordination problem that humanity has ever solved on a global scale
just given the economic incentives at play as AI is becoming a driving
force in innovations of science, technology, and military applications.
If we are fortunate we will have a mild AI armageddon scenario that
compels the global population to enact an international treaty or
moratorium on AI and robotics development.

Yet even with international agreements to curtail new development or
limit deployment of AI, the cat is already out of the bag; in reality we
need to prepare for the worst, with or without future regulations.

Manufacturers and creators may attempt to bake in software-based
controls for restricting AI decision making within constraints, but such
backdoor controls can just as easily be exploited: in particular China
has
\href{https://www.forbes.com/sites/zakdoffman/2025/04/17/china-is-everywhere-your-iphone-android-phone-now-at-risk/}{repeatedly}
\href{https://media.defense.gov/2025/Dec/04/2003834878/-1/-1/0/MALWARE-ANALYSIS-REPORT-BRICKSTORM-BACKDOOR.PDF}{exploited}
such backdoors. AGI now or in the future can certainly exploit such
backdoors or controls more easily than humans can. In short the set of
potential interested and capable parties for conducting a mass hack of
military contractors creating killer drones and android humanoid robots
and chips is huge, while defensive capabilities (e.g.~by formally
verified hardware and software) are severely restricted in adoption in
this capitalistic profit-centric environment.

\begin{quote}
`No regrets,' says Edward Snowden, after 10 years in exile But
whistleblower says 2013 surveillance `child's play' compared to
technology today. -
\href{https://www.theguardian.com/us-news/2023/jun/08/no-regrets-says-edward-snowden-after-10-years-in-exile}{Ewen
MacAskill, The Guardian; June 8th, 2023}
\end{quote}

From a popular meme regarding surveillance technology:

\begin{quote}
\begin{itemize}
\tightlist
\item
  The FBI distributes viruses and keyloggers that mainstream anti-virus
  software are legally not allowed to detect. {[}Magic Lantern, CIPAV,
  Carnivore (DCS1000), Network Investigative Technique{]}
\item
  The NSA has forced Intel, AMD, and chip makers to backdoor their CPUs
  and allow them to access your computer even if it is ``turned off''
  (as long as it has access to electricity) {[}Intel ME, AMD PSP{]}
\item
  The NSA has also forced hardware manufacturers to backdoor their
  `Random Number Generators' to allow them to break RSA encryption
  {[}Dual Elliptic Curve{]}
\item
  American-made electronics transmit radio frequencies which allow the
  FBI and NSA to access your computer even if it's not connected to the
  internet {[}Cottonmouth-I, SURLYSPAWN, ANT/TAO Catalog{]}
\item
  The NSA intercepts computers/laptops/phones purchased online and
  installs malware chips (see: ANT/TAO Catalog) on them before
  delivering them to you {[}TAO team{]}
\item
  If you are a ``Person of Interest'' the NSA is intercepting your mail,
  opening them up, and installing bugs on any electronics you order over
  the internet (e.g.~TVs, microwaves, phones, video games, smart
  refrigerators) {[}ANT/TAO Catalog{]}
\item
  The NSA has already broken Tor networks and can identify you and
  attack you if you use Tor; ISP companies like Verizon, AT\&T, Comcast
  assist with this {[}FOXACID, QUANTUMINSERT, EGOTISTICALGIRAFFE{]}
\item
  Major American tech companies like Microsoft willingly provide the NSA
  advanced warnings of zero-day security vulnerabilities before
  revealing them to the public {[}TAO Team, PRISM, NSA hoarding
  program{]}
\item
  The CIA can infect your car's control systems, which would allow them
  to remotely take it over and drive it {[}Vault 7{]}
\item
  The CIA/MI5 can hack into all smart devices (like TVs and Alexa)
  connected to the internet and spy on you {[}Weeping Angel{]}
\end{itemize}
\end{quote}

Recent Amazon datacenter outages were caused by administrative error
where internal AI was allowed to make changes to production systems
without human oversight. In February 2026 Kim Dotcom
\href{https://x.com/KimDotcom/status/2023165849721536672}{claimed that
Palantir was hacked by an AI agent that gained super-user access},
alleging mass surveillance of world leaders including backdoored
devices, cars, and jets. AI armageddon global scenarios are much more
terrifying when considering all the backdoors that could be exploited by
rogue AGI. Some of these backdoors may even be hijacked to cause
physiological and psychological effects upon nearby people. All of these
backdoors create more surface area for vulnerability making AI
armageddon more certain and more deadly.

\subsection{Problem: Central Bank Digital
Currencies}\label{problem-central-bank-digital-currencies}

Central Bank Digital Currencies (CBDCs) represent perhaps the most
immediate threat to individual financial freedom. As of 2025, 137
countries and currency unions are exploring CBDCs. China's digital yuan
pilot has processed over 7 trillion e-CNY (\textasciitilde\$986 billion)
in transactions. Russia plans Digital Ruble transactions by September
2026. Brazil's Drex CBDC launches in 2026. The European Central Bank's
digital euro is expected between 2026 and 2029. The UAE's Digital Dirham
targets a full launch by late 2026.

Unlike physical cash or even existing bank deposits, CBDCs are
\emph{programmable money}: governments can program expiration dates,
restrict what you can purchase, freeze accounts instantly without
judicial process, implement negative interest rates that force spending,
and surveil every transaction in real time. This is not
speculation--these are documented design features. China's e-CNY already
implements programmable restrictions. The Bank for International
Settlements (BIS) has
\href{https://www.bis.org/speeches/sp220118.htm}{explicitly stated} that
a ``key difference with CBDCs is that the central bank will have
absolute control on the rules and regulations that will determine the
use of that expression of central bank liability, and also, we will have
the technology to enforce that.''

CBDCs are the technological infrastructure for the ``social credit''
model of governance: comply, and your money works; dissent, and your
money stops. This is the ``mark of the beast'' not as metaphor but as
monetary engineering. During the Canadian trucker protests of 2022 the
government froze bank accounts of protesters and donors--with CBDCs such
actions require no cooperation from banks at all.

The antidote is sound money that no central authority controls: gold,
silver, Bitcoin, and well-governed blockchain-based financial systems.
Gno.land can host DeFi applications, stablecoins backed by real assets,
and governance systems for community-managed currencies that operate
outside the CBDC paradigm. The ease of developing complex financial
applications in Gno as compared to Solidity makes it uniquely suited for
the rapid development of these alternatives before the window closes.

Beyond financial sovereignty, our civilization is at risk of a global
catastrophe on par or much worse than the historic burning of the
Library of Alexandria (the largest and most significant library of the
ancient world). Like the burning of the Library of Alexandria, the loss
of our digital knowledge is likely to be accidental, simply because we
did not prepare enough for it. My greatest desire is for Gno.land to
help decentralize the internet archives for preservation and processing
for the benefit of future generations.

\subsubsection{Use Case: Gno.land for Data Archival and
Processing}\label{use-case-gno.land-for-data-archival-and-processing}

Aaron Swartz wrote in his
\href{https://ia800101.us.archive.org/1/items/GuerillaOpenAccessManifesto/Goamjuly2008.pdf}{2008
Guerilla Open Access Manifesto}:

\begin{quote}
``The world's entire scientific and cultural heritage, published over
centuries in books and journals, is increasingly being digitized and
locked up by a handful of private corporations. \ldots{} Forcing
academics to pay money to read the work of their colleagues? Scanning
entire libraries but only allowing the folks at Google to read them?
Providing scientific articles to those at elite universities in the
First World, but not to children in the Global South? It's outrageous
and unacceptable. \ldots{} Those with access to these
resources---students, librarians, scientists--you have been given a
privilege. You get to feed at this banquet of knowledge while the rest
of the world is locked out. But you need not--indeed, morally, you
cannot---keep this privilege for yourselves.''
\end{quote}

There exist a number of archival projects that share books and journals
such as Anna's Archive, Library Genesis, Sci-Hub, UbuWeb and Z-Library.
Take for example Library Genesis:

\begin{quote}
``Library Genesis (shortened to LibGen) is a shadow library project for
file-sharing access to scholarly journal articles, academic and
general-interest books, images, comics, audiobooks, and magazines. The
site enables free access to content that is otherwise paywalled or not
digitized elsewhere.'' -
\href{https://en.wikipedia.org/wiki/Shadow_library}{Wikipedia}
\end{quote}

Some of these projects have or are working toward IPFS integration on
top of BitTorrent; but for the value that global archiving provides not
just today but also for all possible potential future timelines the
number of seeders for many of the archival files
\href{https://zrthstr.github.io/libgen_torrent_cardiography/index.html}{still
hover in the single digits}.

More recently the US Department of Justice released millions of files of
the ``Epstein Files'' many of which have also been made available for
sharing (such as by
\href{https://github.com/yung-megafone/Epstein-Files/tree/main/Torrent\%20Files}{lists
of torrents}.

Some projects such as thewebb.io are developing web tools for analyzing
the Epstein Files and other high profile topics with the assistance of
AI. But you can only know if the AI is being truthful:

\begin{enumerate}
\def\labelenumi{\arabic{enumi}.}
\item
  if you can replicate the response deterministically on another
  computation instance (unless you have direct access to the AI
  hardware, or you are given a complete diagnostic record of internal
  computation, but practically you need repeatable determinism to rule
  out one-off injections of misinformation) and
\item
  if you completely understand the design of the AI and its training
  data set; because otherwise you can't know whether the AI even can
  provide a reasonable answer to the question, or whether it was trained
  with biased data to give the wrong answers.
\end{enumerate}

By making everything deterministic, by committing the Merkle root of
each data collection (whether Library Genesis or the Epstein files) and
also committing the Merkle roots of all the relevant AI code and design
(for deterministic build, but also audits for each release version), the
training corpora, and all other information needed to inspect the
substance of the AI image and also to deterministically build it and use
it; only then can we be given some assurances about the answer. If we
don't know likewise how the answer was generated, then that creates the
opportunity for absolute abuse; and as we have seen with the censorship
of all major Web2.0 platforms this opportunity is taken, such as with
Google when it censors anything contrary to the position of the WHO
regarding infectious diseases.

For this purpose Gno contracts should incentivize more IPFS and torrent
seed nodes for sharing public archive files while keeping the files off
the main blockchain. While such projects surely already exist,
applications and libraries are much easier to develop in a general
purpose language like Go/Gno, and Gno's innovative multi-user awareness
makes it the ideal language for building complex applications composed
of interoperable modules.

While one set of modules on Gno.land can help manage the availability of
data, another set of modules can process the data and build artifacts to
be stored themselves as archival files. Open source tokenized indexers
can create search indices on these archives, even those populated by web
crawlers. By decomposing the search engine and inverting control away
from centralized entities to public accountable micro-services
registered on chain users have the power to participate, verify, and
fund for data gathering, indexing, and processing as they wish. The end
result is a more trustworthy search engine; and for many the ability to
search without being tracked for advertisers.

No longer would we need to depend on a handful of centralized search
engines, and the few internet archives (many thanks to archive.org and
other archives, but they are
\href{https://blog.archive.org/2024/05/28/internet-archive-and-the-wayback-machine-under-ddos-cyber-attack/}{constantly
hit} by
\href{https://blog.archive.org/2024/10/21/internet-archive-services-update-2024-10-21/}{cyberattacks}
and are always under existential risk of being shut down by the powers
that be or even mobs of terrorists or religious extremists. We can
reinvent the data archives and search engines to be more open,
decentralized, accountable/deterministic, and to a large degree more
permissionless and participatory; not to replace the centralized
alternatives, but for the sake of having alternatives that are provably
trustworthy.

Furthermore we should also develop malware exploit detection tools to
scan these files of these archives (especially the PDF files, but for
all types of files) for malware and exploit hacks; anti-virus for public
documents without relying on any central party. Clearly all popular
technology is compromised and we need to get back to basics, even avoid
the usage of PDF files and complex ever-changing web browsers entirely;
but in practice our ability to avoid exploits
\href{https://cryptome.org/2013/12/appelbaum-30c3.pdf}{is almost nil} in
the face of overreaching state actors and private entities. Since no
state seems to be up for the task at hand, we seem to need a
well-organized and well-funded public intelligence agency partially
governed by a chain that is tasked specifically with creating provably
secure (or in the very least finished or near-complete software that
doesn't mutate with the latest trends) foundational open-source software
implementations that all align with ``small. simple. secure.''
philosophy of projects like Tendermint or
\href{https://alpinelinux.org}{Alpine Linux}.

We also have the ability to construct \emph{new} small/simple/secure
composable modules (in Gno and Go) for the new opportunities of today
(such as processing these archives w/ fully accountable deterministic
AI). Once we have attained some level of financial power as a community
we can also start funding the completion of foundational software
projects written in more modern memory-safe languages, formally verify
their safety, and even use AI to help secure them to finality.

Going back to the previous example of using AI to analyze archives,
answers to common queries should be recorded somewhere again with full
accountability of the entire chain of data and logic needed to compute
the query responses such that anybody can easily verify the result for
themselves. Some of these responses could be stored directly on Gno.land
but they could also be stored and indexed in another archive that
specializes in the accountable memoization of common queries.

\subsection{Gno Language}\label{gno-language}

\emph{We could not have arrived at the final design of the Gno language
without the help of the many developers who gave much time of their
lives to contribute to this project and the design of the language;
specifically the NewTendermint Gno team and Gno community, and
especially Manfred, Morgan, Maxwell, Milos, Guilhem, Ray, and Omar.}

\textbf{Gno is the first general purpose \emph{multi-user} programming
language.}

By ``multi-user'' we mean a language where (1) multiple users' code and
data coexist in a shared persistent heap, (2) trust boundaries between
users are enforced by the language itself--not by the operating system,
a hypervisor, or ad-hoc middleware--and (3) inter-user interoperability
uses the same syntax and type system as intra-user code. Prior systems
have achieved subsets of these properties: MUMPS (1966) shares
persistent global state across users but has no language-level trust
boundaries; MOO (1990) and LPC (1989) provide multi-user object
ownership in virtual worlds but are domain-specific MUD languages, not
general purpose; Erlang (1986) excels at concurrency but all processes
are equally trusted with no per-user security model; Solidity (2015) and
Move (2019) handle multi-user smart contracts but lack transparent
inter-module interoperability with a shared heap (as detailed
\hyperref[gno-vs-existing-smart-contract-platforms]{below}). Gno is the
first language to combine all three properties in a general purpose
language based on an existing mainstream language (Go).

Gno is a minimal extension of the Go language for multi-user
programming. Gno allows a massive number of programmers to iteratively
and interactively develop a single shared program such as Gno.land
(making it ideal for interoperable smart contract programming). In other
words, Go is a restricted subset of the Gno language in the single-user
context.

All of our programming languages to date are designed for a single
programmer user. All programming languages make the same assumption that
there is only one user -- the programmer, or program executor user.
Whether written in C, C++, Python, Java, Javascript, or Go, it is
assumed that all of the dependencies of the program are trusted. If
there is a vulnerability in any of the dependencies there is a
vulnerability in the program; it is the job of the programmer or
program/product manager to ensure that the overall program is free of
exploits.

When interacting with programs owned by another user (or process)
various techniques are used such as via IPC APIs often generated by
tools like Protobuf/GRPC; but such tools add extra complexity,
additional surface area for exploits, additional compute complexity, and
do not benefit directly from the language's native rules and
type-checker--especially for inter-process passing of in-memory object
references. In Gno, interacting with another realm's function is just
like interacting with any dependency (import and call).

Automatic persistence of memory objects makes databases, ORMs, and
serialization generally unnecessary. GnoVM at the end of realm
transaction boundaries will ``finalize'' by scanning all new, modified,
and deleted objects and will automatically persist what needs to be
saved along with meta information like the Merkle hash, ref count,
associated realm and unique number of each object. By making inter-user
interoperability native to the language and also by removing the
distinction between volatile and persistent memory, Gno eliminates much
of the accidental complexity that plagues traditional multi-service and
smart-contract development.

The GnoVM is implemented purely in Go and is a Gno AST interpreter. It
is implemented to serve as a reference implementation suitable for
production and tinkering. The GnoVM opcodes are expected to change over
time to balance simplicity, efficiency, and transpilation
targets--perhaps even to hardware. Just as Ethereum was originally
implemented in Python we expect future software to be optimized, whether
implemented in other languages, or possibly transpiled from the existing
GnoVM implementation.

For the latest Gno specs including the inter-realm spec refer to the
\href{https://github.com/gnolang/gno/tree/master/docs/resources}{Gno
docs/resources directory}.

\subsubsection{Gno vs Existing Smart Contract
Platforms}\label{gno-vs-existing-smart-contract-platforms}

Smart contract platforms like Ethereum allow for many users to upload
their application and call other user application logic functions, but
Solidity is not a general purpose programming language and has severe
limitations that make it suboptimal for solving the task at hand.

First, \textbf{Solidity and other existing smart-contract
languages/platforms do not support transparent inter-module (inter-user)
interoperability with the same language rules as for intra-module
dependencies.} That is, an application developer for a smart-contract
cannot simply import and call/use another user application's functions
and types as if they were library dependencies of the same application.
Generally interoperability between different modules are implemented
with extra-language frameworks and libraries on top of an incomplete or
primitive message-passing agent architecture; such interop function
calls generally do not share the same call-stack nor memory space.
Solidity's implicit \passthrough{\lstinline!msg.sender!} changes on
external calls have caused billions of dollars in exploits (reentrancy,
confused-deputy attacks). In Gno the programmer writes
\passthrough{\lstinline!cross!} at each call site that changes the
realm-context, making every trust boundary visible in the source code
and checkable by the compiler--a small syntactic cost for a huge safety
gain. Move (used by Sui and Aptos) improves on Solidity with a linear
type system that prevents reentrancy, but at the cost of a bespoke
language unfamiliar to most developers and without a shared persistent
object graph; inter-module calls in Move still pass values by copy or
move rather than by shared reference.

Second, \textbf{Solidity and other existing smart-contract
languages/platforms do not support the automatic persistence and
Merkle-ization of in-memory (heap) objects and often require custom
serialization logic.} Solidity does not support such a heap at all as
all memory for variables are predeclared in the function and as such is
not object-oriented and does not have a garbage collector or similar
memory-management primitives. WASM-based smart contract systems do not
support automatic persistence of objects without persisting the entire
memory state of the module. This requires a specialized virtual machine
such as the Gno VM which keeps track of every object created, modified,
and deleted. Solidity developers manage storage slots and mappings
manually; Move developers serialize structs into global storage with
explicit \passthrough{\lstinline!move\_to!} and
\passthrough{\lstinline!borrow\_global!} calls. Gno developers write
ordinary Go structs and they persist automatically with
Merkle-ization--a major reduction in complexity for the common case.

The automatic persistence of in-memory objects of the GnoVM is like a
memristor simulator. The advent of AI has created a new market for
memristor-based memory systems where the distinction between RAM
volatile memory and persistent disk storage is removed. Urbit is similar
but is not based on any general purpose programming language. With
memristor-based memory the GnoVM can be further simplified and the
performance of applications can be vastly improved without any changes
to the Gno language specification.

Third, \textbf{Solidity and other existing smart-contract
languages/platforms do not support a shared heap memory space for
objects to be referenced by external-user objects in a uniform manner by
language rules}. Alice cannot simply declare a structure object that
references the structure object persisted in Bob's application and trust
the garbage collector to retain Bob's object for as long as Alice's
object is retained.

Fourth, \textbf{the split between immutable p-packages
(\passthrough{\lstinline!gno.land/p/!}) and mutable r-realms
(\passthrough{\lstinline!gno.land/r/!}) solves a real trust problem that
no other smart contract platform addresses at the package-system level.}
Two mutable realms cannot trustlessly cooperate because either could
upgrade its code, but an immutable p-package can enforce a contract
between them. In Solidity, library contracts can be proxied and
upgraded, undermining any guarantees they were supposed to provide. In
Move, modules are also upgradeable by default. Gno bakes the
immutability guarantee into the package path convention itself.

\textbf{The above differentiating factors of the Gno language allows for
the most succinct expression of a single-user application or multi-user
application composed of independent modules without the extra complexity
from extra-language interop type-checking syntax or frameworks nor of
the extra complexity from any database, ORM, or serialization logic.}

Shared garbage-collection in a shared (multi-user) graph of object
references makes it possible for one's object representing (say) a
propositional statement or idea to be easily referenced by an
alternative statement or idea, or even be extended by reference with
additional commentary, metadata, or even a subreddit-like tree of
discussions. Without a shared garbage collector the task of ensuring
that references still hold over time without becoming dangling pointers
is left up to each inter-application interface at best, requiring custom
logic just to handle garbage collection. WebAssembly (WASM) externref
support in Go has limitations, particularly in how it handles external
memory references. Currently, Rust and Go do not natively support
externref types for function parameters or return values, making it
challenging to pass complex data between Wasm modules and their host
environments effectively.

\begin{quote}
Reference type (aka externref or anyref) is an opaque reference made
available to a WASM module by the host environment. Such references
cannot be forged in the WASM code and can be associated with arbitrary
host data, thus making them a good alternative to ad-hoc handles (e.g.,
numeric ones). References cannot be stored in WASM linear memory; they
are confined to the stack and tables with externref elements.

Rust does not support reference types natively; there is no way to
produce an import / export that has externref as an argument or a return
type. wasm-bindgen patches WASM if externrefs are enabled. This library
strives to accomplish the same goal for generic low-level WASM ABIs
(wasm-bindgen is specialized for browser hosts).

\textbf{externref use cases} Since externrefs are completely opaque from
the module perspective, the only way to use them is to send an externref
back to the host as an argument of an imported function. (Depending on
the function semantics, the call may or may not consume the externref
and may or may not modify the underlying data; this is not reflected by
the WASM function signature.) An externref cannot be dereferenced by the
module, thus, the module cannot directly access or modify the data
behind the reference. Indeed, the module cannot even be sure which kind
of data is being referenced. -
https://docs.rs/externref/latest/externref/
\end{quote}

Even if externref were fully implemented in future specs for Go (or
Rust) such that it could be used as an argument or return type across
modules (still not ideal for type-checking as it is not the underlying
type), this would limit what can be inter-module-referenced to that
which can be held in memory. The Gno Virtual Machine (GnoVM) allows for
inter-user-package (inter-realm) references across the entire persisted
disk store space, and does not require any additional language syntax
such as with the \passthrough{\lstinline!externref!} keyword, and
supports the normal course of type-checking already familiar to Go
developers.

\subsubsection{Interrealm Programming}\label{interrealm-programming}

Gno extends Go with explicit multi-user realm awareness. The key
concepts:

\begin{itemize}
\tightlist
\item
  \textbf{Realm-context}: determines
  \passthrough{\lstinline!runtime.CurrentRealm()!} and
  \passthrough{\lstinline!runtime.PreviousRealm()!}, controlling
  identity and agency.
\item
  \textbf{Realm-storage-context}: determines where new and modified
  objects are persisted during finalization.
\item
  \textbf{Crossing-functions}: declared with
  \passthrough{\lstinline!func fn(cur realm, ...)!} and called with
  \passthrough{\lstinline!fn(cross, ...)!} to explicitly change both
  contexts.
\item
  \textbf{Readonly taint}: values accessed via dot-selectors or
  index-expressions on external realm objects are tainted read-only,
  preventing direct mutation even by logic declared in the same realm.
  This protects against a class of exploits where a realm is tricked
  into mutating its own objects.
\item
  \textbf{Realm-transaction finalization}: at realm boundaries, new
  objects are assigned IDs and persisted, ref-count-zero objects are
  deleted, and Merkle hashes are recomputed.
\end{itemize}

\begin{lstlisting}[language=Go]
// realm /r/bob/bob
package bob

import "gno.land/r/alice/alice" // import external realm package
import "gno.land/p/bob/types"   // import external library package

func Register(cur realm, name string) {
    prof := types.UserProfile{Name: name}
    alice.SetObject(cross, prof) // explicit realm crossing
}
\end{lstlisting}

Non-crossing methods on receivers residing in external realms implicitly
change only the realm-storage-context (not the realm-context), allowing
methods to behave identically whether declared in a p package or realm
package. Panics that cross realm boundaries abort the program;
\passthrough{\lstinline!revive(fn)!} provides software transactional
memory semantics for testing (and eventually production use). Future
builtins include \passthrough{\lstinline!attach()!} for binding unreal
objects to a realm-storage-context, and
\passthrough{\lstinline!safely(cb)!} for clearing realm-contexts to
prevent yielding agency to callees.

For the full specification including the object model
(real/unreal/escaped objects), readonly taint rules, crossing semantics,
and design goals, see the
\href{./resources/gno-interrealm.md}{Interrealm Specification}.

\subsection{Gno.land Blockchain}\label{gno.land-blockchain}

Tendermint solved proof-of-stake by innovating upon classical Byzantine
fault-tolerant consensus algorithms published by Dwork, Lynch, and
Stockmeyer in 1988 (originally funded by Darpa for missile defense
systems) for blockchains. It paved the way for the Cosmos Hub, the first
proof-of-stake IBC hub, and Cosmos the internet of blockchains. Also of
note, when Binance first launched they used the CosmosSDK and
Tendermint.

Gno.land builds upon Tendermint2 and aims to shift the paradigm of
programming languages in general: Gno is the first \emph{multi-user}
programming language, making it a superior smart contracting language as
compared to any existing solution. Thus \textbf{Gno.land is the first
multi-user language-based operating system}. Its ultimate goal is to be
the world's open knowledge base for the next millennium.

\subsubsection{Use Case: Open Programmable Knowledge
Base}\label{use-case-open-programmable-knowledge-base}

Go's simple embedded struct-centric design and the Gno VM's automatic
transactional persistence makes Gno.land not only great for
decentralized financial applications but also makes it uniquely well
suited and designed for permissionless innovation of information-based
applications such as social communication and coordination systems, or
the next Wikipedia or programmable knowledge-base systems. The latter
will be explored here.

Each of the thought statements in the
\hyperref[gno-land-for-mass-awakening]{introduction} can be represented
as a simple Go string, but as in Tractatus we want to allow each of
these thought- statements to be supported by any number of supporting
thought statements, so we need a struct declaration.

\begin{lstlisting}[language=Go]
type Thought struct {
    Statement    string
    Dependencies []*Thought
}
\end{lstlisting}

The above allows for a simple tree structure, but it would be better to
annotate each child node (thought statement) with the type of relation
to the parent node-- for example whether a child node represents an
example, a caveat, a corollary, or supporting evidence and so on.

\begin{lstlisting}[language=Go]
// Option "Denormalized Thought"
type Thought struct {
    Statement   string
    Examples    []*Thought
    Caveats     []*Thought
    Corollaries []*Thought
    Support     []*Thought
}
\end{lstlisting}

Better than a denormalized structure is one where the type of thought
statement association can be arbitrary or fixed depending on the
application.

\begin{lstlisting}[language=Go]
// Option "Normalized Thought"
type Thought struct {
    Text         string
    TypedSupport []*Thought
}

type ThoughtType string // examples, caveats, corollaries, support

type TypedThought struct {
    Type    ThoughtType
    Thought *Thought
}
\end{lstlisting}

\emph{Note on the usage of \passthrough{\lstinline![]\\*Thought!}
slices: in the current implementation of the GnoVM each slice can only
be used by first loading the entire underlying base array. This may be
optimized in the future, however for holding large sets of elements the
programmer should instead use a tree-structure such as the avl.Tree (or
an iavl.Tree).}

Now arises the question of whether counter-arguments should also be
referenced as a child node to the original thought parent node. If we
include counter-arguments in the graph of
\passthrough{\lstinline!*Thought!} objects itself there is the issue of
permissioning who can add counter-arguments to the graph. With the
examples above and with no method declarations a
\passthrough{\lstinline!*Thought!} belonging to one user cannot be
modified by a third party even though the fields of a
\passthrough{\lstinline!Thought!} struct is exposed due to Gno (runtime)
interrealm rules that taint third party reads via direct dot-selectors
\& index-expressions with a readonly-taint that persists even with
(direct selector) access of sub-fields.

The \passthrough{\lstinline!*Thought!} object can however be modified by
another user by calling a declared method. We can extend the
\passthrough{\lstinline!Thought!} struct with additional fields for
authorization or ownership and implement a method such as follows:

\begin{lstlisting}[language=Go]
type Thought struct {
    Owner        account
    Statement    string
    TypedSupport []*Thought
}

func (th *Thought) AddCounterArgument(cth *Thought) {
    caller := runtime.CurrentRealm().Address
    if th.Owner != caller {
        panic("unauthorized")
    }
    th.TypedSupport = append(th.TypedSupport,
        TypedThought{Type: "counter", Thought: cth})
}
\end{lstlisting}

This works but not well--it only works if the owner of the parent node
wants the counter-argument to be registered. Even if counter-arguments
were not registered as an assocation on chain, it is still possible for
any Gno.land state indexer to separately index the reverse association
of reference to the original \passthrough{\lstinline!*Thought!} when it
finds a counter-argument \passthrough{\lstinline!*Thought!} that
references in its struct field the original as a counter-argument. This
reliance on an external indexer shifts trust from the blockchain itself
to the indexer so is not always ideal.

Gno is intended for permissionless iteration and improvement so we will
discuss another way (among many) to manage associations of competing
thought statements; the pair-wise association among competing thought
statements can be registered in another (neutral) external realm that
allows the registration only at least one of the two thought statements
identify the other as a counter-argument. In this case it is not
necessary for a \passthrough{\lstinline!*Thought!} object to be
associated with any owner explicitly (via the
\passthrough{\lstinline!.Owner!} field). Note however that given the Gno
inter-realm specification to make a \passthrough{\lstinline!*Thought!}
object truly immutable even for the owner of the realm in which it
resides it must not expose any mutator functions, or it should have at
least a \passthrough{\lstinline!readonly bool!} field.

We can also add discussion board objects for each thought statement.

\begin{lstlisting}[language=Go]
// Board defines a type for boards.
type Board struct {
    // ID is the unique identifier of the board.
    ID ID
    // Name is the current name of the board.
    Name string
    // Aliases contains a list of alternative names for the board.
    Aliases []string
    // Readonly indicates that the board is readonly.
    Readonly bool
    // Threads contains all board threads.
    Threads PostStorage
    // ThreadsSequence generates sequential ID for new threads.
    ThreadsSequence IdentifierGenerator
    // Permissions enables support for permissioned boards.
    // This type of boards allows managing members with roles and permissions.
    // It also enables the implementation of permissioned execution of board related features.
    Permissions Permissions
    // Creator is the account address that created the board.
    Creator address
    // Meta allows storing board metadata.
    Meta any
    // CreatedAt is the board's creation time.
    CreatedAt time.Time
    // UpdatedAt is the board's update time.
    UpdatedAt time.Time
}

// New creates a new basic non permissioned board.
func New(id ID) *Board {
    return &Board{
        ID:              id,
        Threads:         NewPostStorage(),
        ThreadsSequence: NewIdentifierGenerator(),
        CreatedAt:       time.Now(),
    }
}
\end{lstlisting}

While it is certainly possible to embed a
\passthrough{\lstinline!*Board!} as a field of each
\passthrough{\lstinline!*Thought!}, the current implementation of
\passthrough{\lstinline!*Board!} is only safe from a moderation
perspective when it is permissioned; and so a board tightly coupled to a
\passthrough{\lstinline!*Thought!} may not be ideal depending on the
use-case. Instead we can map an external realm persisted index of
\passthrough{\lstinline!*Thought!} to \passthrough{\lstinline!*Board!}
associations similarly to how counter-thoughts are associated as
mentioned before. In both cases we probably want to add to the
\passthrough{\lstinline!Thought!} struct a globally unique ID like how
\passthrough{\lstinline!Board!} has. \emph{In the future we may provide
a standard function to get a unique identifier for every pointer object
but this has not yet been decided yet.}

Finally, consider for example the numbered sequence of verses of a book
of the bible, or the deep tree of statements in Wittgenstein's
Tractatus. In order to facilitate the efficient forking of such large
lists or graphs of objects it is necessary to avoid the usage of slices.
Even the avl.Tree (as provided in the Gno monorepo under the examples
directory) is not sufficient as it is a mutable tree. However a fork of
the avl.Tree into an immutable tree (or likewise a port of the iavl tree
in the tm2 Tendermint2 directory) can be used with some improvement to
allow for splicing in new elements and deleting existing elements from
the original tree.

So far I have illustrated a way for multiple users to construct their
thought statement graphs independently while also allowing for
counter-arguments to be registered/associated permissionlessly. More
design and exploration is needed to create a fully functional
permissionlessly extensible thought statement graph system; and in the
primordial soup of Gno ecosystem eventually one or more designs will
become dominant in usage by evolution. The reader is encouraged to
explore the above template and measure success by references and by
forks. See also \hyperref[use-case-95-facts]{Use Case: 95 Facts}.

\subsubsection{Use Case: Home Computing
Chains}\label{use-case-home-computing-chains}

The Gno VM is not just useful in the context of public decentralized
blockchains. It is also useful for home computing. Take for example
Email which despite all attempts to replace it still persists in our
lives today as flawed, complex, and outdated as it is. While a realm
that stores mail on the blockchain is not useful unless the data is
encrypted \emph{(and even if it were encrypted it is not a good idea nor
encouraged to store encrypted data on gno.land as encryption keys may
eventually get acquired by hackers and leaked and even persisted on the
blockchain too)}, the Gno VM can run anywhere, even on your private
server hosted at home. In fact, this is what we should do given the
prevalence of surveillance technology such as Google's Gmail which uses
AI to sort your mail and analyze for targeted advertising.

Imagine a black box local GnoVM you run at home. You can have the
\passthrough{\lstinline!/r/home/email!} realm store your emails at home
on your own home server. The same blockchain node logic can run on its
own as a single-validator \emph{home chain} which naturally supports
backups as secondary full nodes, or you can even make your home chain
byzantine fault-tolerant for better uptime.

\begin{enumerate}
\def\labelenumi{\arabic{enumi}.}
\tightlist
\item
  Install in your home GnoVM chain a service plugin: /s/email/indexer,
  not /p/* nor /r/* but /s/* for off-chain service applications.
  \emph{(this prefix is not supported in gno.land but may be in the
  future)}.
\end{enumerate}

\begin{itemize}
\item
  /s/email/indexer reads state upon init, but also registers as a
  listener for notifications from /r/emails.
\item
  When a new email comes in, /r/emails via listeners calls
  \passthrough{\lstinline!/s/email/indexer.AddEmail()!}.
\item
  /s/email/indexer also imports /d/email/indexer which is an off-chain
  daemon component. Here /d/* represents a hypothetical prefix for Gno
  code to be run off-chain with arbitrary Go native functions available
  for import that would otherwise not be possible on gno.land (since a
  blockchain can only support deterministic logic).
\item
  /d/email/indexer can only access /s/email/indexer by a Gno firewall
  system declared with Gno package paths, types, and function/method
  names.
\item
  /s/email/indexer can import any /r/* or /p/* but not any /s/* (like
  Chrome extensions) and its own /d/email/indexer, unless otherwise
  restricted by the Gno firewall system.
\item
  Your mobile device registers an account with your local GnoVM home
  chain. This phone account is restricted to only access
  /s/email/indexer.
\item
  Your phone makes a request to /s/email/indexer. It then asks
  /d/email/indexer which in turn queries the local index and responds
  via /s/email/indexer.
\end{itemize}

Here are some benefits of GnoVM home computing: * Gno.land can be
leveraged to ensure that all software is properly audited. * Software is
expected to become finished and immutable. * All software benefits from
Go/Gno's type-safety and memory-safety. * A unified IPC system
drastically reduces surface area for penetration. * Plugin services and
daemons such as the aforementioned email indexer can be containerized
and restricted from unauthorized access. * Fine-grained
security/firewall rules at the function invocation level. * Byzantine
fault-tolerance comes out of the box for zero downtime. * Inversion of
control with public key cryptography for everything: no more password
management. * Resilience against infrastructure failures: home chains
continue operating during internet outages, power grid disruptions, or
state-level network shutdowns. Combined with mesh networking and local
power (solar, battery), GnoVM home chains can form the backbone of
community-level communication and coordination when centralized
infrastructure fails.

\subsubsection{Use Case: Supply Chain Transparency and Food
Sovereignty}\label{use-case-supply-chain-transparency-and-food-sovereignty}

The fragility of global supply chains has become impossible to ignore.
From semiconductor shortages to food supply disruptions, dependency on
opaque, centralized, and geographically concentrated supply chains puts
populations at the mercy of forces they cannot see, verify, or
influence. The geopolitical tension surrounding Taiwan and TSMC, the
concentration of pharmaceutical manufacturing in China and India, and
the consolidation of food production under a handful of multinational
corporations all represent single points of failure with civilizational
consequences.

Gno contracts can implement transparent, accountable supply chain
tracking where every step from origin to consumer is recorded with
cryptographic accountability. Unlike centralized supply chain solutions
(which are only as trustworthy as the company operating them), Gno-based
supply chain modules are open source, auditable, and composable. A food
provenance realm can import a certification realm which imports a
reputation realm--all with type-safe interoperability and no middleware.

For food sovereignty specifically, local communities can use Gno.land to
coordinate direct farmer-to-consumer networks, manage
community-supported agriculture (CSA) memberships, and maintain
transparent records of growing practices--all without depending on any
centralized platform that could deplatform them or change terms of
service. When combined with home computing chains, local food networks
can operate entirely independently of the global internet if necessary.

\subsubsection{Other Use Cases}\label{other-use-cases}

Gno.land can be used to host any other smart-contract application
supported by Ethereum written in Solidity, such as Defi applications,
name-resolution systems, DAOs and governance applications, etc.

You can explore the various dapps including sample implementations of
ERC equivalents in the
\href{https://github.com/gnolang/gno/tree/master/examples/gno.land}{examples
directory}. \emph{Note that these prototypes have not yet been audited
unless otherwise specified!}

There are also many good use-cases for integrating the GnoVM in client
or server side programming where deterministic gas-limited (for CPU,
disk, memory) execution of logic is desired. In particular Gno and the
Gno VM can be ideal for scripting in the gaming context for multi-player
games. There are many games like Minecraft (with JavaScript) or Second
Life (with LSL) that offer scripting in the game, but they are usually
not both deterministic and gas-counted; but these are the properties you
may need in a multi-player game if you also want to provide
fault-tolerant, accountability, and replayability guarantees (with or
without decentralization or replication).

The main feature of Gno is its multi-user context awareness. Where-ever
you want user provided logic directly interacting with those of another
user Gno solutions will be more succinct, elegant, and easy to develop
and maintain.

\subsubsection{Gno.land Constitution}\label{gno.land-constitution}

See {[}./CONSTITUTION.md{]} for the Gno.land Constitution (draft) and
details of genesis, tokenomics, governance, and more.

\paragraph{Separation of Church and
State}\label{separation-of-church-and-state}

Madison separated church and state in the US Constitution albeit there
is a hint of the Christian spirit by the way in which the constitution
was signed: ``\ldots{} in the Year of the Lord\ldots{}''. All the
founders were Christian including Jefferson and Madison, and in
particular the primary author of the US Constitution James Madison
explicitly separated church from the constitution so as to help promote
the teachings of Jesus as evidenced in his other writings. Likewise
Gno.land besides this whitepaper is independent of any religion by its
constitution, which should only refer to this whitepaper sparingly.

Gno.land will launch with a minimal (living) constitution written and
maintained in English, but also ultimately be supplemented by the
completed GnoVM code and Tendermint2 and Gno.land implementation. Future
implementations of the GnoVM and Gno.land should adhere to the completed
software mentioned above.

Gno.land should not censor speech, even if the speech is wrong. However,
it should ban all porn and try to limit external links to porn sites as
porn is not speech and is dangerous to civilization. Whether hate-speech
is tolerated should be determined by each realm but also by the living
Gno.land constitution and by GovDAO vote to amend the constitution and
laws of Gno.land.

\subsubsection{GnoWeb Browser}\label{gnoweb-browser}

GnoWeb is the server software for Gno.land, a browser within a browser
for viewing realm data.

Instead of requiring realm applications to return HTML, the convention
is to implement a Render() function that returns Markdown. This is to
allow the transition away from the bloated HTML standards and browser
software and realm data to be browsed even from the console. Note that
we don't need HTML XML elements to denote objects: everything in Gno is
already an object. This makes Gno.land more like the original World Wide
Web that conforms to the \emph{Document Object} model (DOM).

GnoWeb does support some custom XML elements for improved layout and
functionality, such as column layout and form submissions that integrate
with browser extensions for transaction signing.

Note that GnoWeb is not yet a general blockchain explorer (e.g.~for
transactions) nor a general purpose Gno.land state explorer. A Gno.land
specific blockchain explorer already exists. GnoWeb can only render
markdown returned from Render() functions, and a general purpose state
explorer is still desired.

Realm code is not precluded from returning HTML or even JSON for custom
browser applications. In the near future the Gno.land node software will
support returning JSON encodings of Gno objects. Thus future alternative
browser applications may provide more interactive rich user experiences
for viewing and mutating Gno.land state without any Markdown
intermediary representation; and perhaps leveraging AI for intelligent
layout and styling for rich interactivity.

\subsection{Future Work}\label{future-work}

\begin{itemize}
\tightlist
\item
  Access control systems.
\item
  Persistent garbage collector
\item
  (and possibly a tx based cyclic garbage collector)
\item
  Name registry for realm and package paths.
\item
  Realm upgrading.
\item
  Realm data browser.
\item
  Deterministic concurrency.
\item
  Joeson parser.
\item
  Gno2.
\item
  Open hardware.
\end{itemize}

\subsection{How to Participate}\label{how-to-participate}

The problems outlined in this manifesto are immense but so is the
opportunity. Every line of Gno code, every realm deployed, every node
operated is a brick in the foundation of a more resilient civilization.
Here is how you can contribute:

\begin{itemize}
\tightlist
\item
  \textbf{Learn Gno.} If you know Go you already know most of Gno. Start
  with the
  \href{https://github.com/gnolang/gno/tree/master/examples/gno.land}{examples
  directory} and the \href{https://gno-by-example.com/}{Gno by Example}
  tutorials.
\item
  \textbf{Write a realm.} Build something useful--a board, a registry, a
  token, a game, a supply chain tracker, a reputation module. Deploy it
  to the testnet. The ecosystem grows one realm at a time.
\item
  \textbf{Run a node.} Validator nodes secure the network. Full nodes
  preserve the state. Home chain nodes give you sovereignty. The more
  nodes the more resilient the network.
\item
  \textbf{Audit and review.} Read other people's realms. Find bugs.
  Suggest improvements. Good code review is as valuable as new code.
\item
  \textbf{Fork.} If you disagree with the direction of Gno.land, fork
  it. The constitution explicitly protects this right. Forkability is
  not a threat to the project--it is the ultimate guarantee of its
  integrity.
\item
  \textbf{Build offline-first.} Develop applications that work on home
  chains and mesh networks. When the internet goes down these
  applications become the infrastructure that communities depend on.
\item
  \textbf{Spread the word.} Share this manifesto. Translate it. Discuss
  it. The ideas here only matter if people know about them.
\end{itemize}

You do not need permission to participate. That is the point.

\subsection{Summary}\label{summary}

The convergence of AI-driven threats, the collapse of dollar hegemony,
the rise of CBDCs and programmable surveillance money, institutional
capture of health and information systems, and the fragility of
centralized infrastructure all point to the same conclusion: we need to
build alternatives now, not after the crisis arrives.

There is only so much that we can do as individuals, and it is not clear
what we can do even as a well organized collective when the incentives
are stacked against the masses. But we need not feel disheartened
because there is much to be built in the domain of open source
foundational tooling for our present and future civilizations. We can
always find joy in constructing a better future brick by brick,
metaphorically speaking.

Gno.land is not merely a blockchain or a programming language. It is an
operating system for human coordination that does not require trust in
any central authority. Every realm is a small act of sovereignty. Every
p-package is a shared agreement that cannot be unilaterally altered.
Every home chain is a household that owns its own data and computation.
Together they form something that no corporation, government, or rogue
AI can easily destroy: a distributed, forkable, accountable foundation
for the next civilization.

\begin{itemize}
\tightlist
\item
  Gno is the next C for multi-user environments.
\item
  Gno.land is the world's new open programmable knowledge base.
\end{itemize}

\subsection{Appendix}\label{appendix}

\subsubsection{Gno.land for Mass
Awakening}\label{gno.land-for-mass-awakening}

\begin{quote}
By the time of Trajan in 117 AD, the denarius was only about 85 percent
silver, down from Augustus's 95 percent. \ldots{} But the real crisis
came after Caracalla, between 258 and 275, in a period of intense civil
war and foreign invasions. The emperors simply abandoned, for all
practical purposes, a silver coinage. By 268 there was only 0.5 percent
silver in the denarius. - Joseph R. Peden, ``Inflation and the Fall of
the Roman Empire''
(\href{https://mises.org/mises-daily/inflation-and-fall-roman-empire}{link})
\end{quote}

Consider the following thought statement tree/graph:

\begin{itemize}
\tightlist
\item
  The Federal Reserve and the fiat dollar is unconstitutional and
  therefore illegitimate.

  \begin{itemize}
  \tightlist
  \item
    The US Constitution Article 1, Section 8, Clause 5 explicitly states
    that ``The Congress shall have Power\ldots{} To coin Money, regulate
    the Value thereof, and of foreign Coin, and fix the Standard of
    Weights and Measures.''
  \item
    The US Constitution Article 1, Section 10, Clause 1 explicitly
    states that ``No state shall \ldots{} coin Money; emit Bills of
    Credit; make any Thing but gold and silver Coin a Tender in Payment
    of Debts.''

    \begin{itemize}
    \tightlist
    \item
      The first ``greenback'' paper dollar issued in 1862 was a bill of
      credit backed by the federal government's promise to pay the
      bearer gold or silver.
    \end{itemize}
  \item
    Neither Article 1, Section 8, Clause 5 nor Article 1, section 10,
    Clause 1 of the US Constitution give the federal government the
    authority to stray from the ``bimetalic'' spirit of the U.S.
    Constitution.
  \item
    Paper fiat dollar bills are not coins.

    \begin{itemize}
    \tightlist
    \item
      Paper fiat dollars started off as bills of credit for deposited
      gold.
    \item
      Paper fiat dollars today are neither bills of credit nor
      gold/silver coins.
    \end{itemize}
  \item
    The Federal Reserve was unconstitutionally ratified in order to
    debase the people's money from the underlying gold and silver.

    \begin{itemize}
    \tightlist
    \item
      The Coinage Act of 1873 also known as the ``Crime of 1873''
      eliminated the standard dollar from the list of coins that the
      U.S. Mint could issue, beginning the demonetization of silver in
      favor of gold with the new ``trade dollar'', and later the
      introduction of the gold dollar coin.
    \item
      The U.S. fiat paper dollar became gold-backed with the passage of
      the Gold Standard Act on March 14, 1900, which established gold as
      the exclusive backing for the country's paper currency. This meant
      that each dollar bill was convertible into a specific amount of
      gold.
    \item
      The Bretton Woods system was an international monetary system
      established in 1944 that set rules for commercial and financial
      relations among the major industrial states pegging the dollar to
      gold.

      \begin{itemize}
      \tightlist
      \item
        It created a system of fixed exchange rates, with the U.S.
        dollar pegged to gold and other currencies pegged to the dollar,
        and established the International Monetary Fund (IMF) and the
        World Bank ``to promote economic stability and growth''.
      \item
        The system ended in 1971 when the U.S. abandoned the gold
        standard, leading to a shift to floating exchange rates.
      \item
        The Bretton Woods system ended in 1971 primarily because the
        U.S. could no longer maintain the dollar's convertibility to
        gold due to rising inflation and a growing balance of payments
        deficit, leading to a loss of confidence in the dollar.
        President Nixon's decision to suspend gold convertibility on
        August 15, 1971, effectively marked the collapse of the system,
        transitioning the world to floating exchange rates.
      \end{itemize}
    \item
      The Coinage Act of 1963 removed silver from dimes and quarters and
      reduced the silver content of half-dollars to 40\%, more or less
      completing the destruction silver coin as money and ending
      bimetallism.

      \begin{itemize}
      \tightlist
      \item
        \href{https://en.wikipedia.org/wiki/Coinage_Act_of_1873}{Wikipedia
        incorrectly states} that bimetalism ended with the Coinage Act
        of 1873, and
        \href{https://en.wikipedia.org/wiki/Talk:Coinage_Act_of_1873}{prevents}
        users from correcting the record.
      \item
        By 1973, the U.S. dollar was fully decoupled from gold,
        transitioning to a fiat currency not backed by physical
        commodities.
      \end{itemize}
    \item
      Hypothesis: JP Morgan intentionally sank the Titanic to murder
      opposition such as Straus and Astor, specifically to debase the
      dollar and to steal the works of Nikola Tesla.

      \begin{itemize}
      \tightlist
      \item
        JP Morgan sank the Titanic to debase the dollar.

        \begin{itemize}
        \tightlist
        \item
          John Jacob Astor IV was the world's richest man; he opposed
          the Treasury and WWI.
        \item
          Isidor Straus was the elected Treasurer of the New York branch
          of the National Citizen's League for the promotion of a Sound
          Banking System who corresponded with the editor of the New
          York Times and
          \href{./images/manifesto/isador_straus.jpeg}{made a call to
          action in and around 10/16/1911} of the public for open
          discussion to prevent the adoption of an unaccountable federal
          reserve act.
        \item
          Less than one year after Isidor Straus declared the call to
          action to the public, the Titanic sunk in April 15, 1912 after
          boarding members invited by JP Morgan who himself dipped out
          of the party at the last minute--knowing that there was an
          engine room fire as logged by logs.

          \begin{itemize}
          \tightlist
          \item
            The Titanic sunk when the engine room exploded as intended;
            rescue ships that trailed the Titanic did not come to rescue
            upon seeing emergency SOS flairs allegedly because the red
            flares were swapped out for white flares that looked like
            fireworks.

            \begin{itemize}
            \tightlist
            \item
              The recent 3D scan of the sunken Titanic reveals an
              outward blowing out of the ship's hull where the engine
              room was, exactly where photographs of the Titanic showed
              fire damage during boarding.

              \begin{itemize}
              \tightlist
              \item
                The OceanGate submersible that later imploded on the way
                to view the Titanic had in its board of directors a
                Rothschild. (See also
                \href{https://ephesus.church\#babylon-in-banking}{Babylon
                in Banking} and
                \href{https://ephesus.church\#the-rothschild-dynasty}{The
                Rothschild Dynasty}.)
              \end{itemize}
            \end{itemize}
          \end{itemize}
        \item
          The year following the sinking of the Titanic saw the
          unconstitutional passage of the Federal Reserve act in
          December 23, 1913.

          \begin{itemize}
          \tightlist
          \item
            The Federal Reserve Act was drafted in secret on Jekyll
            Island and overseen by banking elites including the
            Rothschilds.
          \end{itemize}
        \item
          Less than one year after the passage of the Federal Reserve
          Act began World War I in July 28, 1914; and thus began the
          significant dilution of the dollar via the sale of government
          war bonds.
        \end{itemize}
      \item
        Hypothesis: JPMorgan sank the Titanic to steal the works of
        Nikola Tesla.

        \begin{itemize}
        \tightlist
        \item
          John Jacob Astor IV was Tesla's primary patron.
        \item
          Tesla died on 7 January 1943, at the age of 86, penniless.
        \item
          Two days after his death the FBI ordered the Alien Property
          Custodian to seize Tesla's belongings, even though Tesla was
          an American citizen.
        \item
          Tesla's entire estate from the Hotel New Yorker and other New
          York City hotels was transported to the Manhattan Storage and
          Warehouse Company under the Office of Alien Property (OAP)
          seal.
        \item
          John G. Trump, a professor at M.I.T. and a well-known
          electrical engineer serving as a technical aide to the
          National Defense Research Committee, was called in to analyze
          the Tesla items in OAP custody.

          \begin{itemize}
          \tightlist
          \item
            John G. Trump was the uncle of Donald J. Trump.
          \end{itemize}
        \item
          After a three-day investigation, Trump's report concluded that
          there was nothing which would constitute a hazard in
          unfriendly hands, stating: ``{[}Tesla's{]} thoughts and
          efforts during at least the past 15 years were primarily of a
          speculative, philosophical, and somewhat promotional character
          often concerned with the production and wireless transmission
          of power; but did not include new, sound, workable principles
          or methods for realizing such results''.
        \end{itemize}
      \end{itemize}
    \item
      The silver dollar used to have 0.7734 troy ounces of silver. At
      the price of silver today at \$84 per troy ounce, a silver dollar
      today would be worth \$64, roughly representing about 6 halvings
      of the purchasing power of the dollar, or 1.56\%.
    \item
      The gold dollar used to weigh 1.672 grams of which 90\% was gold.
      At the price of gold today at \$4,590 per troy ounce, a gold
      dollar today would be worth \$244, roughly representing about 8
      halvings of the purchasing power of the dollar, or 0.41\%.
    \item
      Given 0.7734 troy ounces for a silver dollar and 0.053 troy ounces
      for a gold dollar the ratio of the value of silver to gold used to
      be around 1:14.6. Today the ratio is 1:54.6.
    \item
      If the dollar were to be backed by gold and silver as it should,
      the price of gold and silver would be significantly higher.
    \item
      With the adoption of silver solid-state battery technology and the
      greater need for silver in industry the ratio of the value of
      silver to gold will revert to the historical norm or even higher.
    \end{itemize}
  \end{itemize}
\end{itemize}

(See the appendix on \hyperref[the-history-of-u-s-silver-coins]{The
History of U.S. Silver Coins} for more history of U.S. silver-based
money.)

All empires fall after debasing its currency. The Roman Empire took
centuries to debase their currency to 0.5\%. The United States has
achieved the same in about half a century--it has only been 55 years
since President Nixon ended the convertibility of the dollar to gold.
For a modern comparison, the British Empire ended soon after WWI after
debasing the silver shilling starting from the 1920's until 1947 when
silver content was completely eliminated. It is widely understood that
the US empire is in rapid decline and that much of the world will
experience economic struggle as the once mighty petrodollar (and thus by
extension the fiat money of other nations) collapses.

Furthermore silver is becoming increasingly important for industrial
uses, especially with new silver-based solid state batteries entering
the market that are superior to other battery systems, making the return
of silver coin as money inevitable.

The collapse of dollar hegemony is not a future prediction--it is
happening now. The BRICS nations (Brazil, Russia, India, China, South
Africa, and expanding membership) are actively constructing alternatives
to the dollar-based financial system. In 2025 BRICS central banks added
nearly 800 metric tonnes of gold to their reserves, with combined
holdings now exceeding 6,000 tonnes. China and Russia now conduct most
of their bilateral trade in yuan and rubles, bypassing the dollar
entirely. In October 2025 the bloc piloted the ``Unit'', a settlement
instrument backed by a 40\% gold and 60\% BRICS-currency basket,
signaling a structural move toward de-dollarization. Meanwhile BRICS
Pay, a blockchain-based decentralized payment messaging system, is being
developed to circumvent Western-controlled systems like SWIFT.

The implications are profound: as the dollar loses reserve status the US
government's ability to finance its debt through inflation is curtailed,
which means either severe austerity, debt default, or
hyperinflation--none of which end well for ordinary people holding
dollar-denominated assets. This is not a distant scenario; it is the
trajectory we are on. The need for decentralized, sound-money
alternatives and the tools to build them is urgent.

Thus far in human history the powers that be have been able to control
public consciousness and thus human behavior; even gradually and
systematically moving away from the gold and silver coin monetary system
mandated by the US constitution; but we live in a different
post-singularity era of the information age where complete censorship is
all but impossible, and the people have the power to participate in
shaping public consciousness by the development and usage of open source
software, especially blockchain systems that are much more resilient to
censorship than centralized systems.

In the information age it is completely feasible (and eventually
inevitable) that we evolve beyond the age of empires and build a global
community not based on any centralized global governance system or even
based on the might of any military or police force, but instead based on
the development, propagation, and voluntary adoption of protocols of
game-theoretic Nash equilibrium for the world's benefit. But to solve a
civilization level problem first the people need to understand the
problem.

The above tree of thought statements regarding the US dollar is but one
of many that need to be understood by everyone such that we stop being
manipulated for the benefit of others and transcend the cycle of
exploitation. There are so many more such statements of facts that must
be compiled altogether; each need to be similarly fleshed out into their
own subtrees or branches of supporting thought statements. Some
statements will be more relevant to people in one location or
circumstance while others are more globally applicable.

\subsubsection{Open Innovation for Proof-of-Human, Reputation, and
Moderation}\label{open-innovation-for-proof-of-human-reputation-and-moderation}

It is already the case that a significant portion of posts and comments
on social media platforms are AI generated. We desperately need to build
systems to help distinguish between AI and humans. The way proposed by
the powers that be is a universal ``Digital ID'' system but at the same
time we wish not to be subjected to a universal state controlled
identification system as they are likely to be used to increasingly
censor speech and worse--this is the ``mark of the beast'' in the book
of revelation.

To guarantee our safety in the face of potentially tyrannical
governments we need to collectively create a large variety of
alternative protocols and applications built from the ground up, and
simultaneously avoid using centrally mandated solutions. Even with
proof-of-human protocols the AI problem cannot be solved completely
because people often use AI to generate the content that they post
online; however the worst effects of AI can be mitigated if we can
distinguish between content \emph{uploaded} by people vs those uploaded
by AI.

Many WoT (web-of-trust) based reputation systems had been proposed and
researched but they are largely not used outside of large centralized
companies that keep the information to themselves within walled gardens
for their own benefit. The benefit of these walled gardens is that the
information is more restricted (which helps preserve privacy to some
degree) but the downside as mentioned before is that in reality the data
is already being abused to manipulate the masses for profit. Therefore
the benefit of open WoT reputation systems is that the information
asymmetry in favor of large corporations is nullified.

The obvious downside is that privacy is much reduced unless the
reputation system is designed to preserve it. There are ways of
implementing limited scale privacy preserving reputation systems using
FL (federated learning) and HE or FHE (fully homomorphic encryption)
techniques that have been published or is being implemented. In
particular FHE based systems are interesting because they may be more
quantum resistant.

See also: * \href{https://eprint.iacr.org/2024/506.pdf}{``A
Decentralized Federated Learning using Reputation''} - ``federated
learning (FL) and fully homomorphic encryption (FHE) can be components
of a quantum-safe solution for privacy preserving reputation
computations.'' *
\href{https://www.semanticscholar.org/reader/ac0990f7a063a02222982f798be4bfa206728df2}{Privacy
Preserving PageRank Algorithm By Using Secure Multi-Party Computation} *
\href{https://eprint.iacr.org/2018/917.pdf}{Secure multiparty PageRank
algorithm for collaborative fraud detection} *
\href{https://github.com/hariom2509/universal-fhevm-sdk}{a minimal FHE
EVM implementation for ratings}

Reputation scoring is only one way to filter, restrict access, or
moderate content. Any group or entity (such as a university, local
community, or corporation) can publish a list of pseudonymous accounts
associated with that group, and on a higher level a set of whitelisted
groups can be constructed to create a metagroup for the aforementioned
purposes.

In the near future we may see or need competing (franchise or
independent) storefronts that allow anyone to show up in person to
register a pseudonymous crypto account in a privacy preserving way; such
franchises may partner with any number of auditors that may set up
cameras or other equipment with attestable hardware of the franchise
business's operations, allowing permissionless auditability of these
in-person registrations. Such attestations can allow qualified or
unqualified people to be whitelisted for some applications, and this can
all happen without any government mandated digital ID system.

Other types of components to be implemented include: * Hierarchical
chain of accountability systems (e.g.~large corporation) * Oracles for
Web2.0 accounts * Oracles for citizenship or residency * Reputation
scoring based on source of funds * Deposit or purchase based systems *
Viral invitation based systems
(\href{https://gist.github.com/jaekwon/410a552ce9363d49790ea444a3c00c99}{prototype
WIP})

These are some of the types of components that we need in order to have
a well functioning blockchain for coordination when other systems fail
in the AI armageddon. In Gno these components can be composed by simply
importing the component's realm path
(e.g.~\passthrough{\lstinline!import "gno.land/r/my\_org/member\_attestation"!})
and using its exposed functions just as if it were any other dependency
library.

To illustrate a concrete example consider the management of discussion
boards on Gno.land
(\href{https://gno.land/r/gnoland/boards2/v1:OpenDiscussions}{example}):

\begin{quote}
Realm Package \textgreater{} Board Object \textgreater{} Post Object
\textgreater{} Comment Object
\end{quote}

The realm, board, post, and comment may each have independent access
control policies to determine who gets to add more content. Now Consider
a board that has no controls: such as with a small fixed deposit per
post. Repeated violations by the same account could be detected at the
mempool level or on-chain, but a sophisticated attacker will have many
accounts that are not obviously associated with each other and may have
the capital to take down the whole board because of its insufficient
moderation ability.

With insufficient moderation Gno.land ends up being harmful to users,
just as 4chan.org is dangerous to browse (albeit for the same reason
that it is useful due to its relative lack of censorship); to create a
platform where alternatives can be implemented but with a slight
improvement to eliminate porn and gore without risking freedom-of-speech
is an unstated objective of my drafts of the genesis Gno.land
constitution). \emph{Epstein was allegedly involved in the creation of
4chan.org/pol/; perhaps it's not surprising that it is the way it is}.

With authoritarian moderation GovDAO ends up dictating what policies are
needed to participate (or who gets to enforce them) and wielding such
power necessarily corrupts it such that it potentially becomes the mark
of the beast, not just on Gno.land. Gno.land GovDAO should not require
any specific whitelisting policy or set of policies. Even if a realm is
being abused ideally its legitimate usage should not be affected and
state only purged of specific offending material following strict
procedure; but in practice this may be difficult to enforce especially
when the app is designed such that it blends legitimate and illegitimate
use cases together.

The partial solution to the authoritarian moderation problem is for
neither GovDAO nor any DAO to dictate a solution method but to allow the
realm controller (if any) to make changes to state within the scope of
what it is allowed (as defined by the spirit of the code but possibly
also with a manifest/constitutional document for that realm) to inject a
new permissions interface implementation (moderation policy) that it
sees fit. By giving the realm controller freedom of control and by using
objective metrics to enforce some level of neutrality the free market
has the opportunity to collaboratively, competitively, and
compositionally construct many models and objects for access control.

When the chain's governance fails to seek the middle ground for
moderation purposes it is ultimately the responsibility of the users to
abandon the chain and to seek forks or alternatives. For this purpose
the Gno.land constitution allows for a minority (more than one third) of
GovDAO to establish an official ``qualified'' fork of Gno.land; and in
practice anyone can create an unqualified fork as well with new
governance (or none at all) if abiding by license terms becomes of
secondary import.

This is where I see the greatest need for more Gno code innovation for
Gno.land: a growing collection of control, reputation, and membership
systems that can be composed by anyone such that there is no universal
mandated system but rather a variety solutions that work well enough and
gives users freedom of choice and freedom of association. The spirit of
the constitution and genesis of the chain is one that preserves the
freedom of choice, speech, and association of the user even while
keeping the chain generally free of certain offensive material. If not,
a subset of GovDAO, or in the worst case the users themselves should
fork the chain.

\subsubsection{On E-Voting}\label{on-e-voting}

While we should be innovating on blockchain-based voting systems to
prepare against the worst case scenarios, I cannot yet suggest that
blockchains be used for general elections; not until e-voting systems
can satisfy some prerequisites. From
\href{https://pdfs.semanticscholar.org/6e22/77fbe289333946cd21bfbbeeb0ea7e3c947b.pdf}{one
survey paper on e-Voting systems}:

\begin{quote}
Functional requirements define the desired end functions and features
required by a system. The functional requirements can be directly
observed and tested.

\begin{itemize}
\tightlist
\item
  \textbf{Robustness}: Any dishonest party cannot disrupt elections.
\item
  \textbf{Fairness}: No partial tally is revealed.
\item
  \textbf{Verifiability}: The election results cannot be falsified.
  There are two types of verifiability:

  \begin{itemize}
  \tightlist
  \item
    \textbf{Individual verifiability}: The voter can verify whether
    their vote is included in the final tally.
  \item
    \textbf{Universal verifiability}: All valid votes are included in
    the final tally and this is publicly verifiable.
  \end{itemize}
\item
  \textbf{Soundness, completeness and correctness}: The final tally
  included all valid ballots.
\item
  \textbf{Eligibility}: Unqualified voters are not allowed to vote.
\item
  \textbf{Dispute-freeness}: Any party can publicly verify whether the
  participant follows the protocol at any phase of the election.
\item
  \textbf{Transparency}: Maximise transparency in the vote casting, vote
  storing and vote counting process while preserving the secrecy of the
  ballots.
\item
  \textbf{Accuracy}: The system is errorless and valid votes must be
  correctly recorded and counted. This properties can be retained by
  universal verifiability.
\item
  \textbf{Accountability}: If the vote verification process fails, the
  voter can prove that he has voted and at the same time preserving vote
  secrecy.
\item
  \textbf{Practicality}: The implementation of requirements and
  assumptions should be able to adapt to large-scale elections.
\item
  \textbf{Scalability}: The proposed e-voting scheme should be versatile
  in terms of computation, communication and storage.
\end{itemize}

A security requirement is the required security functionality that
ensures that the scheme satisfies different security properties to solve
a specific security problem or to eliminate potential vulnerabilities.
Security requirements serve as fundamental security functionality for a
system. Therefore, instead of constructing a custom security approach
for every system, the standard security requirements allow researchers
and developers to reuse the definitions of security controls and best
practices.

\begin{itemize}
\tightlist
\item
  \textbf{Privacy and vote secrecy}: The cast votes are anonymous to any
  party.
\item
  \textbf{Double-voting prevention}, unicity and unreusability: Eligible
  voters cannot vote more than once.
\item
  \textbf{Receipt-freeness}: The voter cannot attain any information
  that can be used to prove how he voted for any party.
\item
  \textbf{Coercion-resistance}: Coercers cannot insist that voters vote
  in a certain way and the voter cannot prove his vote to the
  information buyer.
\item
  \textbf{Anonymity}: The identity of the voter remains anonymous.
\item
  \textbf{Authentication}: Only eligible voters were allowed to vote.
\item
  \textbf{Integrity}: The system can detect the dishonest party that
  modifies the election results.
\item
  \textbf{Unlinkability}: The voter and his vote cannot be linked
\end{itemize}
\end{quote}

The paper also lists ``post-quantum'' e-voting schemes, but notes that
it is a new area of research (as of 2022) and that few have implemented
such a system. I like survey papers like these because they usually have
good information in them. I don't agree with all of the conclusions in
the paper but it's worth reviewing along with others.

But does e-voting really make sense? Take for example
``Coercion-resistance''. Some blockchain-based e-voting systems claim to
be coercion-resistant but assume that the adversary (coercer, or vote
buyer) cannot learn of the voter's private key. When voting in person in
a booth theoretically one cannot be coerced if noone besides the voter
knows who they voted for (though this is not the case). But in reality
desperate people will \emph{willingly} give up their private keys and
devices for money. Therefore arguably in all cases electronic voting
that can be performed anywhere are naturally more vulnerable to coercion
and vote-buying.

The only fool-proof voting system in my opinion is one where the voter
is marked with dye also referred to as ``election ink''. On the other
hand, a mark on the right hand or forehead is the hallmark of the ``mark
of the beast'' in the Book of Revelation. One could avoid the right hand
(Ha! some countries use a finger of the left hand) but the key
distinguishing factor of the mark of the beast is that everyone is
required to participate whether rich or poor.

The Book of Luke starts with the circumstances of Jesus's birth in
Bethlehem as it relates to Caesar's universal census registration
mandate. The Book of Revelation ends with a new heaven and new earth and
New Jerusalem right after the opening of the Lamb's book of life where
names and deeds are written of those who are saved. Any system that
tries to duplicate any universal identification system for the purpose
of local governance voting, UBI, or proof-of-person even to distinguish
AI bots from humans has the potential to become a mandated oppressive
solution. The key message of the Book of Revelation and the life of
Jesus as written is against such universal mandatory registries. In
short, any information (identifier) that can be used against someone
\emph{will} be used against them especially in AI armageddon scenarios.

Can Gno.land evolve such that no set of controls is mandated? This is
less of a problem for blockchain systems that are designed primarily for
monetary transactions of fungible tokens, or for blockchain systems that
do not aim to control for any kind of offending material (see
\href{https://github.com/tendermint/atom_one/blob/master/README.md\#how-to-immunize-against-the-mark-of-the-beast}{one
potential solution for monetary systems like Cosmos} that defers
controls to each zone). Yes, but contingent on GovDAO's good judgement
on enforcing a middle ground between insufficient moderation and
authoritarian moderation. Not as the day to day arbiter (except in the
beginning) but as the final court of justice of a system of controls
designed to enforce the spirit of the constitution. If not, GovDAO
itself should fork, and failing that, users should fork the chain or
seek alternatives.

\subsubsection{On Manufactured
Discontent}\label{on-manufactured-discontent}

\begin{quote}
``At the heart of the story are Sigmund Freud's daughter Anna and his
nephew Edward Bernays who had invented the profession of public
relations. Their ideas were \textbf{used by the US government, big
business and the CIA to develop techniques to manage and control the
minds of the American people}. Those in power believed that the only way
to make democracy work and create a stable society was to repress the
savage barbarism that lurked just under the surface of normal American
life.'' - Adam Curtis, BBC Four, 2002.
(\href{https://github.com/jaekwon/ephesus/blob/main/files/century-of-the-self-transcript.pdf}{transcript},
\href{https://www.youtube.com/watch?v=caicn3VpHTo}{youtube})
\end{quote}

\subsubsection{On Twitter Censorship}\label{on-twitter-censorship}

Much can be said about all platforms such as Google, Reddit,
Meta/Instagram, Wikipedia and so on but I will only comment here about
Twitter, also known as ``X''. \emph{Interesting note: the actual name
change happened on Tisha B'av, as I was there when it happened and I
remember thinking it was an ominous date for Elon with the ``beast
armor'' profile to rename Twitter to the letter of a ``mark'' (since
Tisha B'av is a commemoration of a number of disasters in Jewish
history, primarily the destruction of both Solomon's Temple by the
Neo-Babylonian Empire and the Second Temple by the Roman Empire in
Jerusalem); but furthermore news reports retroactively have tried to
change history by claiming that it happened on another date.}

By some reason I got to participate in Twitter's Community Notes program
so I tried it. The first note I was obligated to rate was of a community
note that claimed Elon Musk's wealth did not come from tax payers. Fair
enough, but a little suspicious that this is the first one shown.

A few weeks later after I got used to the mechanics of the system I
noticed a community note by the alias of ``Hilarious Wind Sandgrouse''
(everyone is given a pen name to preserve privacy--mine is ``Charitable
Star Piculet'') that seemed interesting by DerrickEvans4WV about
\href{https://x.com/i/communitynotes/t/2009097879106015609}{Senator
Machaela Cavanaugh who the original poster said ``vandalized'' a
Declaration of Independence exhibit at the state capital building}. The
community note by ``Hilarious Wind Sandgrouse'' only expressed her point
of view and that she said the object was prohibited; but when I read the
\href{https://nebraskaexaminer.com/2026/01/07/nebraska-state-lawmaker-removes-part-of-conservative-pragerus-founders-museum/}{article
referenced by Sandgrouse himself} it was made clear that the objects
were never prohibited, so therefore she was ``vandalizing''. So I looked
further into his account and discovered that (a) Sandgrouse is a
\href{./images/manifesto/twitter_hws_profile.png}{prolific note
contributor}, and (b) Sandgrouse had been
\href{./images/manifesto/twitter_hws_example.jpeg}{spamming biased
community notes} with absurd notes regarding a recent anti-ICE incident
where he claimed that there is no evidence of race motive or
anti-immigration sentiment even though he also wrote that a bullet round
had ``Anti-ICE'' written on it! So did what any responsible person would
do--I went over to each of those tweets to comment (only for community
note contributors to see)
\href{./images/manifesto/twitter_hws_callout.jpeg}{that Sandgrouse is
wildly biased with citations} and while I wasn't sure what would happen,
I didn't expect Twitter to
\href{./images/manifesto/twitter_cm1.jpeg}{suspend by community note
writing privileges} claiming that my recent notes have been rated
``unhelpful'' even though \href{./images/manifesto/twitter_cm2.jpeg}{it
isn't the case}.

So there you have it, Twitter/X censors, and the ``Community Notes''
won't save us either. For one, it doesn't make it easy to point out
obvious bias from people as I did because instead of displaying a
discussion tree as in Reddit you are limited to writing more community
notes (that won't get shown) just to express a concern to the other
community notes contributors (who not only write but also must rank
notes); and besides, Twitter controls when community note contributors
get to vote on notes and by the time a correction is displayed the tweet
may have already peaked in virality. At this point one can almost assume
that Twitter does exactly that.

Besides the above, there's a lot more to be said about Twitter but I'll
leave a few highlights.

\begin{itemize}
\item
  One tweet by @unusual\_whales regarding FTX and tokenized shares from
  Jan 15th, 2023 has \href{./images/manifesto/twitter_ftx.jpeg}{many
  community notes} that claim that Twitter had been suppressing the
  votes by removing likes and retweets.
\item
  My followers count has been stuck at around \textasciitilde26.5K users
  for years, whereas it used to grow exponentially. This occurred around
  the time that I started getting vocal about Covid19 vaccines and
  mandates (and it was revealed that military agencies like DARPA was
  involved in suppressing accounts).
\item
  I have experienced seeing my tweet go viral, with 13 retweets within a
  few minutes, only to see the retweet count drop down to 2 within
  seconds.
\item
  My account used to show time and time again ``shadowban'' status as
  detected by shadowban detection services; but it is clear that my
  account is now suppressed via means that are not so easily detectable.
\item
  This is why I created a new account in the first place, and my
  \href{https://x.com/jaesustein}{old account} I can no longer access,
  and my email address is no longer associated with it.
\end{itemize}

\newpage
\part{Interrealm Specification}
\label{interrealm-specification}
\setcounter{section}{2}
\setcounter{subsection}{0}

\subsection{Introduction}\label{introduction-2}

All modern popular programming languages are designed for a single
programmer user. Programming languages support the importing of program
libraries natively for components of the single user's program, but this
does not hold true for interacting with components of another user's
(other) program. Gno is an extension of the Go language for multi-user
programming. Gno allows a massive number of programmers to iteratively
and interactively develop a single shared program such as Gno.land.

The added dimension of the program domain means the language should be
extended to best express the complexities of programming in the
inter-realm (inter-user) domain. In other words, Go is a restricted
subset of the Gno language in the single-user context. (In this analogy
client requests for Go web servers don't count as they run outside the
server program).

\subsubsection{Interrealm Programming
Context}\label{interrealm-programming-context}

Gno.land supports three types of packages: - \textbf{Realms
(\passthrough{\lstinline!/r/!})}: Stateful user applications (smart
contracts) that maintain persistent state between transactions -
\textbf{Pure Packages (\passthrough{\lstinline!/p/!})}: Stateless
libraries that provide reusable functionality - \textbf{Ephemeral
Packages (\passthrough{\lstinline!/e/!})}: Temporary code execution with
MsgRun which allows a custom main() function to be run instead of a
single function call as with MsgExec.

For an overview of the different package types in Gno
(\passthrough{\lstinline!/p/!}, \passthrough{\lstinline!/r/!}, and
\passthrough{\lstinline!/e/!}), see
\href{../builders/anatomy-of-a-gno-package.md}{Anatomy of a Gno
Package}.

Interrealm programming refers to the ability of one realm to call
functions in another realm. This can occur between: - Regular realms
(\passthrough{\lstinline!/r/!}) calling other regular realms via MsgExec
and MsgRun. - Ephemeral realms (\passthrough{\lstinline!/e/!}) calling
regular realms via MsgRun (like main.go)

The key concept is that code executing in one realm context can interact
with and call functions in other realms while leveraging the language
syntax rules of Go, enabling complex multi-user interactions while
maintaining clear boundaries and permissions.

\subsubsection{Realm-Context and
Realm-Storage-Context}\label{realm-context-and-realm-storage-context}

All logic in Gno executes under two contexts that together govern
identity and persistence:

\textbf{Realm-context} determines
\passthrough{\lstinline!runtime.CurrentRealm()!} and
\passthrough{\lstinline!runtime.PreviousRealm()!}. It controls identity
and agency: who is the current actor and who called them. The
realm-context has an associated Gno address from which native coins can
be sent and received. It changes only on explicit cross-calls
(\passthrough{\lstinline!fn(cross, ...)!}).

\textbf{Realm-storage-context} determines where new and modified objects
are persisted during realm-transaction finalization. It changes on
explicit cross-calls \emph{and} on implicit borrow-crosses (calling a
non-crossing method on a real object residing in a different realm). The
realm-storage-context is not directly accessible at runtime and has no
associated address.

After an explicit cross-call, both contexts refer to the same realm.
They diverge when calling a non-crossing method of a real object
residing in a different realm --- the realm-storage-context shifts to
the receiver's realm while the realm-context stays the same.

\begin{small}
\begin{longtable}{|p{4.5cm}|c|c|c|c|}
\hline
\textbf{Call type} & \textbf{Realm-ctx?} & \textbf{Storage-ctx?} & \textbf{Boundary?} & \textbf{Finalizes?} \\
\hline
\endhead
\texttt{fn(cross, ...)} to same realm & Yes* & No & Yes & Yes \\
\hline
\texttt{fn(cross, ...)} to different realm & Yes & Yes & Yes & Yes \\
\hline
\texttt{fn(nil, ...)} (non-crossing-call) & No & No & No & No \\
\hline
Non-crossing method, real receiver in same realm & No & No & No & No \\
\hline
Non-crossing method, real receiver in different realm & No & Yes & Yes & Yes \\
\hline
Non-crossing method, unreal receiver & No & No & No & No \\
\hline
Non-crossing function & No & No & No & No \\
\hline
\end{longtable}
\end{small}

* \passthrough{\lstinline!runtime.CurrentRealm()!} returns the same
realm, but \passthrough{\lstinline!runtime.PreviousRealm()!} shifts ---
what was current becomes previous. See \hyperref[realm-boundaries]{Realm
Boundaries} for definitions of boundary and finalization.

\subsubsection{Design Goals}\label{design-goals}

\textbf{Caveat: The interrealm specification does not secure
applications against arbitrary code execution. It is important for realm
logic (and even p package logic) to ensure that arbitrary (variable)
functions (and similarly arbitrary interface methods) are not provided
by malicious callers; such arbitrary functions and methods whether
crossing (or non-crossing) will inherit the previous realm (or both
current and previous realms) and could abuse these realm-contexts.} It
does not make sense for any realm user to cross-call an arbitrary
function or method as it loses agency while being marked as the
responsible caller by the callee's runtime previous realm. This problem
is worse when calling a non-crossing function or method. It can be
reasonable when such variable functions or interface values are
restricted in other ways such as by whitelisting by a DAO upon careful
inspection of every such variable function or interface value (both its
type declaration as well as its state).

P package code should behave the same even when copied verbatim in a
realm package; and likewise non-crossing code should behave the same
when copied verbatim from one realm to another. Otherwise there will be
lots of security related bugs from user error.

Realm crossing with respect to
\passthrough{\lstinline!runtime.CurrentRealm()!} and
\passthrough{\lstinline!runtime.PreviousRealm()!} must be explicit and
warrants type-checking; because a crossing-function of a realm should be
able to call another crossing-function of the same realm without
necessarily crossing (changing the realm-context). Sometimes the
previous realm and current realm must be the same realm, such as when a
realm consumes a service that it offers to external realms and users.

For Go developers familiar with blockchain VMs like the EVM: in
Solidity, calling another contract implicitly shifts
\passthrough{\lstinline!msg.sender!}, making it easy to introduce
reentrancy bugs or misattribute the caller. Gno's
\passthrough{\lstinline!cross!} keyword makes every realm-context
transition visible in source code and verifiable by the compiler,
eliminating this class of bugs by construction.

Where a real object resides should not matter too much, as it is often
difficult to predict. Thus the realm-context as returned by
\passthrough{\lstinline!runtime.PreviousRealm()!} and
\passthrough{\lstinline!runtime.CurrentRealm()!} should not change with
non-crossing method calls, and the realm-storage-context should be
determined for non-crossing methods only by the realm-storage of the
receiver. The realm-storage of a receiver should only matter for when
elements reside in external realm-storage and direct dot-selector or
index-expression access of sub-elements are desired of the
aforementioned element.

A method should be able to modify the receiver and associated objects of
the same realm-storage as that of the receiver.

A method should be able to create new objects that reside in the same
realm by association in order to maintain storage realm consistency and
encapsulation and reduce fragmentation.

It is difficult to migrate an object from one realm to another even when
its ref-count is 1; such an object may be deep with many descendants of
ref-count 1 and so performance is unpredictable.

Code declared in p packages (or declared in ``immutable'' realm
packages) can help different realms enforce contracts trustlessly, even
those that involve the caller's current realm. Otherwise two mutable
(upgradeable) realms cannot export trust unto the chain because
functions declared in those two realms can be upgraded. This is
analogous to Go interfaces but stronger: a Go interface guarantees
method signatures, while a Gno p-package guarantees \emph{behavior} ---
the implementation is immutable on-chain, so both parties can trust the
logic without trusting each other.

Both \passthrough{\lstinline!fn(cross, ...)!} and
\passthrough{\lstinline!func fn(cur realm, ...)\{...\}!} may become
special syntax in future Gno versions.

\begin{lstlisting}[language=Go]
// realm /r/alice/alice
package alice

var object any

func SetObject(cur realm, obj any) {
    object = obj
}
\end{lstlisting}

\begin{lstlisting}[language=Go]
// package /p/bob/types
package types

type UserProfile struct {
    Name string
    ...
}
\end{lstlisting}

\begin{lstlisting}[language=Go]
// realm /r/bob/bob
package bob

import "gno.land/r/alice/alice" // import external realm package
import "gno.land/p/bob/types"   // import external library package

func Register(cur realm, name string) {
    prof := types.UserProfile{Name: name}
    alice.SetObject(cross, prof)
}
\end{lstlisting}

The Gno language is extended to support a
\passthrough{\lstinline!context.Context!}-like argument to denote the
current realm-context of a Gno function. This allows a user realm
function to call itself safely as if it were being called by an external
user, and helps avoid a class of security issues that would otherwise
exist.

\begin{lstlisting}[language=Go]
// realm /r/alice/mail

func SendMail(cur realm, text string) {
    if text == "" {
        // runtime.PreviousRealm() is preserved for recursive call.
        SendMail(nil, "<empty>")
    }
    caller := runtime.PreviousRealm()
    if inBlacklist(caller) {
        // runtime.PreviousRealm() becomes self; message from self to self.
        SendMail(cross, fmt.Sprintf("blacklisted caller %v blocked", caller))
    } else {
        // sendMailPrivate not exposed to external callers.
        sendMailPrivate(text)
    }
}
\end{lstlisting}

\subsection{Object Model and Realm
Storage}\label{object-model-and-realm-storage}

Unlike other blockchain platforms where developers must manually
serialize state into key-value stores, Gno persists ordinary Go structs
and slices automatically. A Go developer can write familiar data
structures --- trees, linked lists, maps of structs --- and they survive
across transactions with no explicit marshaling. The runtime handles
Merkle-ization and garbage collection transparently.

Every object in Gno is persisted on disk with additional metadata
including the object ID and an optional OwnerID (if persisted with a
ref-count of exactly 1). The object ID is only set at the end of a
realm-transaction during realm-transaction finalization (more on that
later). A GnoVM transaction is composed of one or many scoped (stacked)
realm-transactions.

\begin{lstlisting}[language=Go]
type ObjectInfo struct {
    ID       ObjectID  // set if real.
    Hash     ValueHash `json:",omitempty"` // zero if dirty.
    OwnerID  ObjectID  `json:",omitempty"` // parent in the ownership tree.
    ModTime  uint64    // time last updated.
    RefCount int       // for persistence. deleted/gc'd if 0.

    // Object has multiple references (refcount > 1) and is persisted separately
    IsEscaped bool `json:",omitempty"` // hash in iavl.
    ...
}
\end{lstlisting}

When an object is persisted during realm-transaction finalization the
object becomes ``real'' (as in it is really persisted in the virtual
machine state) and is said to ``reside'' in the realm; and otherwise is
considered ``unreal''. New objects instantiated during a transaction are
always unreal; and during finalization such objects are either discarded
(transaction-level garbage collected) or become persisted and real.

Unreal (new) objects that become referenced by a real (persisted) object
at runtime will get their OwnerID set to the parent object's storage
realm, but will not yet have its object ID set before realm-transaction
finalization. Subsequent references at runtime of such an unreal object
by real objects residing in other realms do not override the OwnerID
initially set, so during realm-transaction finalization it ends up
residing in the first realm it became associated with (referenced from).
Unreal objects that become persisted but were never directly referenced
by any real object during runtime will only get their OwnerID set to the
realm of the first real ancestor.

Real objects with ref-count of 1 have their hash included in the sole
parent object's serialized byte form, thus an object tree of only
ref-count 1 descendants are Merkle-hashed completely.

When a real or unreal object ends up with a ref-count of 2 or greater
during realm-transaction finalization its OwnerID is set to zero and the
object is considered to have ``escaped''. When such a real object is
persisted with ref-count of 2 or greater it is forever considered
escaped even if its ref-count in later transactions is reduced to 1.
Escaped real objects do not have their hash included in the parent
objects' serialized byte form but instead are Merkle-ized separately in
an iavl tree of escaped object hashes (keyed by the escaped object's ID)
for each realm package. (This is implemented as a stub but not yet
implemented for the initial release of Gno.land.)

\textbf{A real object can only be directly mutated through dot-selectors
and index-expressions if the object resides in the same realm as the
current realm-storage-context. Unreal objects can always be directly
mutated if their elements are directly exposed.}

Exposed values accessed through dot-selectors and index-expressions from
external realm logic are tainted read-only. For the full rules see
\hyperref[readonly-taint-specification]{Readonly Taint Specification}.

\subsection{Crossing-Functions and
Crossing-Methods}\label{crossing-functions-and-crossing-methods}

Realm crossing occurs when a crossing function (declared as
\passthrough{\lstinline!func fn(cur realm, ...)\{...\}!}) is called with
the Gno \passthrough{\lstinline!fn(cross, ...)!} syntax.

\begin{lstlisting}[language=Go]
package main
import "gno.land/r/alice/realm1"

func main() {
    bread := realm1.MakeBread(cross, "flour", "water")
}
\end{lstlisting}

(In Linux/Unix operating systems user processes can cross-call into the
kernel by calling special syscall functions, but user processes cannot
directly cross-call into other users' processes. This makes the GnoVM a
more complete multi-user operating system than traditional operating
systems.)

Besides explicit realm crossing via the
\passthrough{\lstinline!fn(cross, ...)!} Gno syntax, implicit realm
crossing occurs when calling a method of a receiver object stored in an
external realm. Implicitly crossing into (borrowing) a receiver object's
storage realm allows the method to directly modify the receiver as well
as all other objects directly reachable from the receiver stored in the
same realm as the receiver. Unlike explicit crosses, implicit crosses do
not shift or otherwise affect the current realm context;
\passthrough{\lstinline!runtime.CurrentRealm()!} does not change unless
a method is called like
\passthrough{\lstinline!receiver.Method(cross, args...)!}.

Realms hold objects in residence and they also have a Gno address to
send and receive coins from. Coins can only be spent from the current
realm context.

\subsubsection{Definition}\label{definition}

A crossing-function or crossing-method is that which is declared in a
realm and has as its first argument \passthrough{\lstinline!cur realm!}.
The \passthrough{\lstinline!cur realm!} argument must appear as the
first argument of a crossing-function or crossing-method's argument
parameter list. To prevent confusion it is illegal to use anywhere else,
and cannot be used in p packages.

The current realm-context and realm-storage-context changes when a
crossing-function or crossing-method is called with the
\passthrough{\lstinline!cross!} keyword in the first argument as in
\passthrough{\lstinline!fn(cross, ...)!}. Such a call is called a
``cross-call'' or ``crossing-call''.

\begin{lstlisting}[language=Go]
package main
import "gno.land/r/alice/extrealm"

func MyMakeBread(cur realm, ingredients ...any) { ... }

func main(cur realm) {
    MyMakeBread(cross, "flour", "water") // ok -- cross into self.
    extrealm.MakeBread(cross, "flour", "water") // ok -- cross into extrealm
}
\end{lstlisting}

When a crossing-function or crossing-method is called with
\passthrough{\lstinline!nil!} as the first argument instead of
\passthrough{\lstinline!cross!} it is called a non-crossing-call; and no
realm-context nor realm-storage-context change takes place.

\begin{lstlisting}[language=Go]
package main
import "gno.land/r/alice/extrealm"

func MyMakeBread(cur realm, ingredients ...any) { ... }

func main(cur realm) {
    MyMakeBread(nil, "flour", "water") // ok -- non-crossing.
    extrealm.MakeBread(nil, "flour", "water") // invalid -- external realm function
}
\end{lstlisting}

To prevent confusion a non-crossing-call of a crossing-function or
crossing-method declared in a realm different than that of the caller's
realm-context and realm-storage-context will result in either a
type-check error; or a runtime error if the crossing-function or
crossing-method is variable.

\subsubsection{Realm-Context Rules}\label{realm-context-rules}

All functions in Gno execute under a realm context as determined by the
call stack. Objects that reside in a realm can only be modified if the
realm context matches.

A function declared in p packages when called:

\begin{itemize}
\tightlist
\item
  inherits the last realm for package declared functions and closures.
\item
  inherits the last realm when a method is called on unreal receiver.
\item
  implicitly crosses to the receiver's resident realm when a method of
  the receiver is called. The receiver's realm is also called the
  ``borrow realm''.
\end{itemize}

A function declared in a realm package when called:

\begin{itemize}
\tightlist
\item
  explicitly crosses to the realm in which the function is declared if
  the function is declared as
  \passthrough{\lstinline!func fn(cur realm, ...)\{...\}!} (with
  \passthrough{\lstinline!cur realm!} as the first argument). The new
  realm is called the ``current realm''.
\item
  otherwise follows the same rules as for p packages.
\end{itemize}

\passthrough{\lstinline!runtime.CurrentRealm()!} returns the current
realm-context that was last cross-called to.
\passthrough{\lstinline!runtime.PreviousRealm()!} returns the
realm-context cross-called to before the last cross-call. All
cross-calls are explicit with the \passthrough{\lstinline!cross!}
keyword, as well as non-crossing-calls of crossing-functions and
crossing-methods with \passthrough{\lstinline!nil!} instead of
\passthrough{\lstinline!cross!}.

A crossing function declared in the same realm package as the callee may
be called like \passthrough{\lstinline!fn(cross, ...)!} or
\passthrough{\lstinline!fn(cur, ...)!}. When called with
\passthrough{\lstinline!fn(cur, ...)!} there will be no realm crossing,
but when called like \passthrough{\lstinline!fn(cross, ...)!} there is
technically a realm crossing and the current realm and previous realm
returned are the same.

The current realm and previous realm do not depend on any implicit
crossing to the receiver's borrowed/storage realm even if the borrowed
realm is the last realm of the call stack. In other words
\passthrough{\lstinline!runtime.CurrentRealm()!} may differ from the
Machine's active storage realm (internally
\passthrough{\lstinline!m.Realm!}, i.e.~the borrow realm) when a method
is called on a receiver residing in a foreign realm.

\subsubsection{Implicit Realm-Storage
Borrowing}\label{implicit-realm-storage-borrowing}

Besides (explicit) realm-context changes via the
\passthrough{\lstinline!fn(cross, ...)!} cross-call syntax, implicit
realm-storage-context changes occur when calling a non-crossing method
of a receiver object residing in different realm-storage. The
realm-storage-context ``borrows'' to the receiver's realm, which allows
the method to directly modify its receiver and any objects reachable
from it that reside in the same realm-storage. The realm-context does
not change (so \passthrough{\lstinline!runtime.CurrentRealm()!} and
\passthrough{\lstinline!runtime.PreviousRealm()!} are unaffected; the
agency of the caller remains the same). However, objects reachable from
the receiver but residing in a \emph{different} realm-storage cannot be
modified --- the \passthrough{\lstinline!DidUpdate()!} guard in the
runtime enforces that only objects belonging to the borrowed realm can
be mutated.

This borrowing behavior allows non-crossing methods of receivers to
behave the same whether declared in a realm package or p package: p
package code copied over to a realm package, or realm package code
copied over to another realm, have the exact same behavior. Crossing
methods of a realm package would still behave differently when copied
over to another realm, as crossing-methods always change the
realm-context and realm-storage-context to the declared realm.

\subsubsection{Crossing-Method
Semantics}\label{crossing-method-semantics}

When a method is a crossing-method called as
\passthrough{\lstinline!receiver.Method(cross, args...)!}, both the
realm-context and realm-storage-context change to that of the realm
package in which the type/method is declared (which is not necessarily
the same as where the receiver resides). Such a crossing method-call
cannot directly modify the real receiver if it happens to reside in an
external realm that differs from where the type and methods are
declared; but it can modify any unreal receiver or unreal reachable
objects. As mentioned previously a non-crossing-call of a
crossing-method will fail during type-checking or at runtime if the
receiver resides in an external realm-storage.

Calls of methods on receivers residing in realms different from the
current realm must \emph{not} be called like
\passthrough{\lstinline!fn(cross, ...)!} if the method is not a crossing
function itself, and vice versa. Or it could be said that implicit
crossing is not real realm crossing. (When you sign a document with
someone else's pen it is still your signature; signature:pen ::
current:borrowed)

A crossing method declared in a realm cannot modify the receiver if the
object resides in a different realm. However not all methods are
required to be crossing methods, and crossing methods may still read the
state of the receiver (and in general anything reachable is readable).

\begin{lstlisting}[language=Go]
// Type and methods declared in /r/alice/tokens
package tokens

type Token struct {
    Owner   string
    Balance int
}

// Non-crossing method: borrows receiver's storage realm.
func (t *Token) Debit(amount int) {
    t.Balance -= amount
}

// Crossing method: realm-context and realm-storage-context
// shift to /r/alice/tokens (where the type is declared).
func (t *Token) Transfer(cur realm, to string, amount int) {
    t.Balance -= amount // FAILS if receiver resides outside /r/alice/tokens
}
\end{lstlisting}

\begin{lstlisting}[language=Go]
// /r/bob/bob stores a Token
package bob

import "gno.land/r/alice/tokens"

var myToken *tokens.Token // persisted in /r/bob/bob

func Spend(_ realm) {
    // Non-crossing call: storage-context borrows to /r/bob/bob
    // (receiver's realm). Debit() CAN modify myToken.
    myToken.Debit(10) // ok: borrowed storage matches receiver

    // Crossing call: storage-context shifts to /r/alice/tokens
    // (where the type is declared), not /r/bob/bob (where myToken resides).
    // Transfer() CANNOT directly modify myToken.
    myToken.Transfer(cross, "g1...", 10) // fails: receiver in external realm
}
\end{lstlisting}

New unreal objects reachable from the borrowed realm (or current realm
if there was no method call that borrowed) become persisted in the
borrowed realm (or current realm) upon finalization of the foreign
object's method (or function). (When you put an unlabeled photo in
someone else's scrapbook the photo now belongs to the other person.) In
the future the \passthrough{\lstinline!attach()!} function will prevent
a new unreal object from being taken.

For how crossing rules apply to MsgCall, MsgRun, and package
initialization, see \hyperref[message-types-and-testing]{Message Types
and Testing}.

\subsection{Realm Boundaries}\label{realm-boundaries}

A realm boundary is defined as a change in realm in the call frame stack
from one realm to another, whether explicitly crossed with
\passthrough{\lstinline!fn(cross, ...)!} or implicitly borrow-crossed
into a different receiver's storage realm. A realm may cross into itself
with an explicit cross-call.

When a crossing-function or crossing-method is cross-called it shifts
the ``current'' runtime realm-context to the ``previous'' runtime
realm-context such that
\passthrough{\lstinline!runtime.PreviousRealm()!} returns what used to
be returned with \passthrough{\lstinline!runtime.CurrentRealm()!} before
the realm boundary. The current realm-storage-context is always set to
that of realm-context after cross-calling.

For which call types create boundaries, see the
\hyperref[realm-context-and-realm-storage-context]{summary table} above.
Every crossing-call creates a realm boundary even when there is no
resulting change in realm-context or realm-storage-context. A
non-crossing-call of a crossing-function or crossing-method
(\passthrough{\lstinline!fn(nil, ...)!}) never creates a realm boundary.

\subsection{Realm-Transaction
Finalization}\label{realm-transaction-finalization}

Realm-transaction finalization occurs when returning from a realm
boundary. When returning from a cross-call (with
\passthrough{\lstinline!cross!}) realm-transaction finalization will
occur even with no change of realm-context or realm-storage-context.
Realm-transaction finalization does NOT occur when returning from a
non-crossing-call of a method of an unreal receiver or a real receiver
that resides in the same realm-storage-context as that of the caller.

During realm-transaction finalization all new reachable objects are
assigned object IDs and stored in the current realm-storage-context; and
ref-count-zero objects deleted (full ``disk-persistent cycle GC'' will
come after launch); and any modified ref-counts and new Merkle hash root
computed.

\subsection{Readonly Taint
Specification}\label{readonly-taint-specification}

Go's language rules for value access through dot-selectors \&
index-expressions are the same within the same realm, but exposed values
through dot-selector \& index-expressions are tainted read-only when
performed by an external realm.

The readonly taint prevents the direct modification of real objects by
any logic, even from logic declared in the same realm as that of the
object's storage-realm.

A realm cannot directly modify another realm's objects without calling a
function that gives permission for the modification to occur.

For example \passthrough{\lstinline!externalrealm.Foo!} is a
dot-selector expression on an external object (package) so the value is
tainted with the \passthrough{\lstinline!N\_Readonly!} attribute.

The same is true for \passthrough{\lstinline!externalobject.FieldA!}
where \passthrough{\lstinline!externalobject!} resides in an external
realm.

The same is true for \passthrough{\lstinline!externalobject[0]!}: direct
index expressions also taint the resulting value with the
\passthrough{\lstinline!N\_Readonly!} attribute.

The same is true for
\passthrough{\lstinline!externalobject.FieldA.FieldB[0]!}: the readonly
taint persists for any subsequent direct access, so even if FieldA or
FieldB resided in the caller's own realm-context or realm-storage the
result is tainted readonly.

A Gno package's global variables even when exposed
(e.g.~\passthrough{\lstinline!package realm1; var MyGlobal int = 1!})
are safe from external manipulation
(e.g.~\passthrough{\lstinline!import "xxx/realm1"; realm1.MyGlobal = 2!})
by the readonly taint when accessed directly by dot-selector or
index-expression from external realm logic; and also by a separate
\passthrough{\lstinline!DidUpdate()!} guard when accessed by other means
such as by return value of a function and the return value is real and
external.

A function or method's arguments and return values retain and pass
through any readonly taint from caller to callee. Even if a realm's
function (or method) returns an untainted real object, the runtime guard
in \passthrough{\lstinline!DidUpdate()!} prevents it from being modified
by an external realm-storage-context.

For a realm (user) to manipulate an untainted object residing in an
external realm, a function (or method) can be declared in the external
realm which references and modifies the aforementioned untainted object
directly (by a name declared outside of the scope of said function or
method). Or, the function can take in as argument an untainted real
object returned by another function.

Besides protecting against writing by direct access, the readonly taint
also helps prevent a class of security issue where a realm may be
tricked into modifying something that it otherwise would not want to
modify. Since the readonly taint prohibits mutations even from logic
declared in the same realm, it protects realms against mutating its own
object that it doesn't intend to: such as when a realm's real object is
passed as an argument to a mutator function where the object happens to
match the type of the argument.

Objects returned from functions or methods are not readonly tainted. So
if
\passthrough{\lstinline!func (eo object) GetA() any \{ return eo.FieldA \}!}
then \passthrough{\lstinline!externalobject.GetA()!} returns an object
that is not tainted assuming eo.FieldA was not otherwise tainted. While
the parent object \passthrough{\lstinline!eo!} is still protected from
direct modification by external realm logic, the returned object from
\passthrough{\lstinline!GetA()!} can be passed as an argument to logic
declared in the residing realm of \passthrough{\lstinline!eo.FieldA!}
for direct mutation.

Whether or not an object is readonly tainted it can always be mutated by
a method declared on the receiver.

\begin{lstlisting}[language=Go]
// /r/alice

var blacklist []string

func GetBlacklist() []string {
    return blacklist
}

func FilterList(cur realm, testlist []string) { // blanks out blacklist items from testlist
    for i, item := range testlist {
        if contains(blacklist, item) {
            testlist[i] = ""
        }
    }
}
\end{lstlisting}

This is a toy example, but you can see that the intent of
\passthrough{\lstinline!FilterList()!} is to modify an externally
provided slice; yet if you call
\passthrough{\lstinline!alice.FilterList(cross, alice.GetBlacklist())!}
you can trick alice into modifying its own blacklist--the result is that
alice.BlackList becomes full of blank values.

With the readonly taint \passthrough{\lstinline!var Blacklist []string!}
solves the problem for you; that is, /r/bob cannot successfully call
\passthrough{\lstinline!alice.FilterList(cross, alice.Blacklist)!}
because \passthrough{\lstinline!alice.Blacklist!} is readonly tainted
for bob.

The problem remains if alice implements
\passthrough{\lstinline!func GetBlacklist() []string \{ return Blacklist \}!}
since then /r/bob can call
\passthrough{\lstinline!alice.FilterList(cross, alice.GetBlacklist())!}
and the argument is not readonly tainted.

Future versions of Gno may also expose a new modifier keyword
\passthrough{\lstinline!readonly!} to allow for return values of
functions to be tainted as readonly. Then with
\passthrough{\lstinline!func GetBlacklist() readonly []string!} the
return value would be readonly tainted for both bob and alice.

\subsection{\texorpdfstring{\texttt{panic()} and
\texttt{revive(fn)}}{panic() and revive(fn)}}\label{panic-and-revivefn}

\passthrough{\lstinline!panic()!} behaves the same within the same realm
boundary, but when a panic crosses a realm boundary (as defined in
\hyperref[realm-boundaries]{Realm Boundaries}) the Machine aborts the
program. This is because in a multi-user environment it isn't safe to
let the caller recover from realm panics that often leave the state in
an invalid state.

This would be sufficient, but we also want to write our tests to be able
to detect such aborts and make assertions. For this reason Gno provides
the \passthrough{\lstinline!revive(fn)!} builtin.

\begin{lstlisting}[language=Go]
abort := revive(func() {
    fn := func(_ realm) {
        panic("cross-realm panic")
    }
    fn(cross)
})
abort == "cross-realm panic"
\end{lstlisting}

\passthrough{\lstinline!revive(fn)!} will execute `fn' and return the
exception that crossed a realm boundary during finalization.

This is only enabled in testing mode (for now), behavior is only
partially implemented. In the future
\passthrough{\lstinline!revive(fn)!} will be available for non-testing
code, and the behavior will change such that
\passthrough{\lstinline!fn()!} is run in transactional (cache-wrapped)
memory context and any mutations discarded if and only if there was an
abort.

TL;DR: \passthrough{\lstinline!revive(fn)!} is Gno's builtin for STM
(software transactional memory).

\subsection{\texorpdfstring{\texttt{attach()}}{attach()}}\label{attach}

In future releases of Gno the \passthrough{\lstinline!attach()!}
function can be used to associate unreal objects to the current
realm-storage-context before being passed into a function declared in an
external realm package, or into a method of a real receiver residing in
an external realm-context.

\subsection{\texorpdfstring{\texttt{safely(cb\ func())}}{safely(cb func())}}\label{safelycb-func}

In future releases of Gno the
\passthrough{\lstinline!safely(cb func())!} function may be used to
clear the current and previous realm-context as well as any
realm-storage-context such that no matter what
\passthrough{\lstinline!cb func()!} does the caller does not yield
agency to the callee.

For now this can be simulated by implementing an (immutable
non-upgradeable) realm crossing-function that cross-calls into itself
once more before calling the callback function.

\subsection{Guidelines}\label{guidelines}

P package code cannot contain crossing functions. P package code also
cannot import R realm packages. But code can call named crossing
functions e.g. those passed in as parameters.

You must declare a public realm function to be a crossing function if it
is intended to be called by end users, because users cannot MsgCall
non-crossing functions (for safety/consistency) or p package functions
(there's no point).

Utility functions that are a common sequence of non-crossing logic can
be offered in realm packages as non-crossing functions. These can also
import and use other realm utility non-crossing functions; whereas p
packages cannot import realm packages at all. And convenience/utility
functions that are being staged before publishing as permanent p code
should also reside in upgradeable realms.

Generally you want your methods to be non-crossing. Because they should
work for everyone. They are functions that are pre-bound to an object,
and that object is like a quasi-realm in itself, that could possibly
reside and migrate to other realms. This is consistent with any p code
copied over to r realms; none of those methods would be crossing, and
behavior would be the same; stored in any realm, mostly non-crossing
methods that anyone can call. Why is a quasi-realm self-encapsulated
object in need to modify the realm in which it is declared, by crossing?
That's intrusive, but sometimes desired.

You can always cross-call a method from a non-crossing method if you
need it.

\subsection{Message Types and Testing}\label{message-types-and-testing}

\subsubsection{MsgCall}\label{msgcall}

MsgCall may only call crossing functions. This is to prevent potential
confusion for non-sophisticated users. Non-crossing calls of
non-crossing functions of other realms is still possible with MsgRun.

\begin{lstlisting}[language=Go]
// PKGPATH: gno.land/r/test/test

func Public(_ realm) {

    // Returns (
    //     addr:<origin_caller>,
    //     pkgpath:""
    // ) == testing.NewUserRealm(origin_caller)
    runtime.PreviousRealm()

    // Returns (
    //     addr:chain.PackageAddress("gno.land/r/test/test"),
    //     pkgpath:"gno.land/r/test/test"
    // ) == testing.NewCodeRealm("gno.land/r/test/test")
    runtime.CurrentRealm()

    // Call a crossing function of same realm with crossing
    AnotherPublic(cross)

    // Call a crossing function of same realm without crossing
    AnotherPublic(cur)
}

func AnotherPublic(_ realm) {
    ...
}
\end{lstlisting}

\subsubsection{MsgRun}\label{msgrun}

\begin{lstlisting}[language=Go]
// PKGPATH: gno.land/e/g1user/run

import "gno.land/r/realmA"

func main() {
    // Before main() is called there is an implicit
    // crossing from UserRealm(g1user) to
    // CodeRealm(gno.land/e/g1user/run).

    // Returns (
    //     addr:g1user,
    //     pkgpath:""
    // ) == testing.NewUserRealm(g1user)
    runtime.PreviousRealm()

    // Returns (
    //     addr:g1user,
    //     pkgpath:"gno.land/e/g1user/run"
    // ) == testing.NewCodeRealm("gno.land/e/g1user/run")
    runtime.CurrentRealm()

    realmA.PublicNoncrossing()
    realmA.PublicCrossing(cross)
}
\end{lstlisting}

Realm addresses are derived from package paths via
\passthrough{\lstinline!chain.PackageAddress()!} (defined in
\passthrough{\lstinline!gnovm/stdlibs/chain/address.gno!}):

\begin{itemize}
\tightlist
\item
  \passthrough{\lstinline!chain.PackageAddress("gno.land/r/name123/realm")!}
  --- bech32 from hash(path)
\item
  \passthrough{\lstinline!chain.PackageAddress("gno.land/e/g1user/run")!}
  --- bech32 substring ``g1user''
\end{itemize}

Therefore in the MsgRun file's \passthrough{\lstinline!init()!} function
the previous realm and current realm have different pkgpaths (the origin
caller always has empty pkgpath) but the address is the same.

\subsubsection{MsgAddPackage}\label{msgaddpackage}

A realm package's initialization (including
\passthrough{\lstinline!init()!} calls) executes with current
realm-context of itself.
\passthrough{\lstinline!runtime.PreviousRealm()!} refers to the package
deployer both in global var decls and inside
\passthrough{\lstinline!init()!} functions. After that the package
deployer is no longer provided, so packages need to remember the
deployer in the initialization phase if needed.

\begin{lstlisting}[language=Go]
// PKGPATH: gno.land/r/test/test

func init() {
    // Returns (
    //     addr:<origin_deployer>,
    //     pkgpath:""
    // ) == testing.NewUserRealm(origin_deployer)
    // Inside init() and global var decls
    // are the only time runtime.PreviousRealm()
    // returns the deployer of the package.
    // Save it here or lose it forever.
    runtime.PreviousRealm()

    // Returns (
    //     addr:chain.PackageAddress("gno.land/r/test/test"),
    //     pkgpath:"gno.land/r/test/test"
    // ) == testing.NewCodeRealm("gno.land/r/test/test")
    runtime.CurrentRealm()
}

// Same as in init().
var _ = runtime.PreviousRealm()
\end{lstlisting}

\begin{lstlisting}[language=Go]
// PKGPATH: gno.land/e/g1user/run

func init() {
    // Returns (
    //     addr:g1user,
    //     pkgpath:""
    // ) == testing.NewUserRealm(g1user)
    runtime.PreviousRealm()

    // Returns (
    //     addr:g1user,
    //     pkgpath:"gno.land/e/g1user/run"
    // ) == testing.NewCodeRealm("gno.land/e/g1user/run")
    runtime.CurrentRealm()
}
\end{lstlisting}

The same applies for pure package (\passthrough{\lstinline!/p/!})
initialization. During initialization and tests,
\passthrough{\lstinline!runtime.CurrentRealm()!} can return a package
path that starts with ``/p/''. This is because the package is
technically still mutable during its initialization phase. After
initialization, pure packages become immutable and cannot maintain
state.

\subsubsection{Testing overrides with
stdlibs/testing}\label{testing-overrides-with-stdlibstesting}

The
\passthrough{\lstinline!gnovm/tests/stdlibs/testing/context\_testing.gno!}
file provides functions for overriding frame details from Gno test code.

\passthrough{\lstinline!testing.SetRealm(testing.NewUserRealm("g1user"))!}
is identical to
\passthrough{\lstinline!testing.SetOriginCaller("g1user")!}. Both will
override the Gno frame to make it appear as if the current frame is the
end user signing with a hardware signer. Both will also set
\passthrough{\lstinline!ExecContext.OriginCaller!} to that user. One of
these will become deprecated.

\paragraph{\texorpdfstring{Gno test cases with \texttt{\_test.gno} like
\texttt{TestFoo(t\ *testing.T)}}{Gno test cases with \_test.gno like TestFoo(t *testing.T)}}\label{gno-test-cases-with-_test.gno-like-testfoot-testing.t}

\begin{lstlisting}[language=Go]
// PKGPATH: gno.land/r/user/myrealm
package myrealm

import (
    "chain/runtime"
    "testing"
)

func TestFoo(t *testing.T) {
    // At first OriginCaller is not set.

    // Override the OriginCaller.
    testing.SetRealm(testing.NewUserRealm("g1user"))

    // Identical behavior:
    testing.SetOriginCaller("g1user")

    // This panics now: seeking beyond the overridden origin frame:
    // runtime.PreviousRealm()

    // Simulate g1user cross-calling Public().
    // Produce a new frame to override
    func() {
        testing.SetRealm(testing.NewCodeRealm("gno.land/r/user/myrealm"))

        runtime.PreviousRealm() // "g1user", ""
        runtime.CurrentRealm()  // bech32(hash("gno.land/r/user/myrealm")), "gno.land/r/user/myrealm"

        Public(...) // already in "gno.land/r/user/myrealm"
    }()

    // The following is identical to the above,
    // but not possible in p packages which
    // cannot import realms.
    Public(cross, ...)
}
\end{lstlisting}

\paragraph{\texorpdfstring{Gno filetest cases with
\texttt{\_filetest.gno}}{Gno filetest cases with \_filetest.gno}}\label{gno-filetest-cases-with-_filetest.gno}

\begin{lstlisting}[language=Go]
// PKGPATH: gno.land/r/test/test
package test

import (
    "chain/runtime"
    "testing"

    "gno.land/r/user/myrealm"
)

func init() {
    // XXX Frame not found, there is no deployer for filetests.
    runtime.PreviousRealm()

    // Returns (
    //     addr:chain.PackageAddress("gno.land/r/test/test")
    //     pkgpath:"gno.land/r/test/test"
    // ) == testing.NewCodeRealm("gno.land/r/test/test")
    runtime.CurrentRealm()
}

func main() {
    // There is assumed to be in "frame -1"
    // a crossing from UserRealm(g1user) to
    // CodeRealm("gno.land/r/test/test") before
    // main() is called, so crossing() here
    // is redundant.

    // Returns (
    //     addr:g1user,
    //     pkgpath:""
    // ) == testing.NewUserRealm(g1user)
    runtime.PreviousRealm()

    // Returns (
    //     addr:g1user,
    //     pkgpath:"gno.land/r/test/test"
    // ) == testing.NewCodeRealm("gno.land/r/test/test")
    runtime.CurrentRealm()

    // gno.land/r/test/test cross-calling
    // gno.land/r/user/myrealm:
    myrealm.Public(cross, ...)
}

// Output:
// XXX
\end{lstlisting}

\subsection{Implementation}\label{implementation}

Implementation for \passthrough{\lstinline!runtime.CurrentRealm()!} and
\passthrough{\lstinline!runtime.PreviousRealm()!} are defined in
\passthrough{\lstinline!gnovm/stdlibs/chain/runtime/native.gno!} and
related files in the directory, while overrides for testing are defined
in
\passthrough{\lstinline!gnovm/tests/stdlibs/testing/context\_testing.gno!}.
All stdlibs functions are available unless overridden by the latter.

\subsection{Proposed Changes}\label{proposed-changes}

\passthrough{\lstinline!testing.SetOriginCaller()!} may be deprecated in
favor of
\passthrough{\lstinline!testing.SetRealm(testing.NewOriginRealm(user))!}.

\passthrough{\lstinline!testing.NewCodeRealm(path)!} may be renamed to
\passthrough{\lstinline!testing.NewPackageRealm(path)!}.

\end{document}
